\begin{chapterpage}{Probabilitas}
  \chaptertitle{Probabilitas}
  \label{probability}
  \label{ch_probability}
  \chaptersection{basicsOfProbability}
  \chaptersection{conditionalProbabilitySection}
  \chaptersection{smallPop}
  \chaptersection{randomVariablesSection}
  \chaptersection{contDist}
\end{chapterpage}
\renewcommand{\chapterfolder}{ch_probability}

\index{probability|(}

\chapterintro{Probabilitas merupakan landasan statistika,
  dan Anda mungkin \mbox{sudah}
  mengenal banyak gagasan yang disajikan dalam bab ini.
  Namun, formalisasi konsep probabilitas kemungkinan besar masih
  baru bagi sebagian besar pembaca. \\

  \noindent%
  Walaupun bab ini memberikan landasan teoretis
  bagi gagasan-gagasan pada bab-bab selanjutnya dan membuka jalan
  menuju pemahaman yang lebih mendalam,
  penguasaan konsep-konsep yang diperkenalkan di bab ini
  tidak diperlukan untuk menerapkan
  metode yang diperkenalkan di bagian lain buku ini.}

%  This chapter provides a theoretical foundation for
%  the ideas introduced in later chapters.
%  However, this chapter is not strictly required to
%  understand or apply the methods introduced in the
%  rest of this book.}




\section{Mendefinisikan probabilitas}
\label{basicsOfProbability}

Statistika didasarkan pada probabilitas,
dan walaupun probabilitas tidak diperlukan untuk teknik-teknik
terapan dalam buku ini, pemahaman probabilitas dapat membantu Anda
memahami metode secara lebih mendalam dan membangun landasan yang lebih baik
untuk mata kuliah selanjutnya.


\subsection{Contoh pengantar}

Sebelum membahas gagasan teknis, mari kita lihat
beberapa contoh dasar yang mungkin terasa lebih akrab.

\begin{examplewrap}
\begin{nexample}{“Dadu”, bentuk tunggal dari dice, adalah kubus dengan enam sisi bernomor \resp{1}, \resp{2}, \resp{3}, \resp{4}, \resp{5}, dan \resp{6}. Berapa peluang memperoleh \resp{1} saat melempar dadu?}\label{probOf1}
Jika dadunya adil, peluang memperoleh \resp{1} sama dengan peluang memperoleh bilangan lain. Karena ada enam hasil, peluangnya adalah 1 banding 6 atau, secara ekuivalen, $1/6$.
\end{nexample}
\end{examplewrap}

\begin{examplewrap}
\begin{nexample}{Berapa peluang memperoleh \resp{1} atau \resp{2} pada lemparan berikutnya?}\label{probOf1Or2}
\resp{1} dan \resp{2} merupakan dua dari enam hasil yang mungkin dan sama-sama berpeluang, sehingga peluang memperoleh salah satu dari kedua hasil ini adalah $2/6 = 1/3$.
\end{nexample}
\end{examplewrap}

\begin{examplewrap}
\begin{nexample}{Berapa peluang memperoleh \resp{1}, \resp{2}, \resp{3}, \resp{4}, \resp{5}, atau \resp{6} pada lemparan berikutnya?}\label{probOf123456}
100\%. Hasilnya pasti salah satu dari bilangan tersebut.
\end{nexample}
\end{examplewrap}

\begin{examplewrap}
\begin{nexample}{Berapa peluang tidak memperoleh \resp{2}?}\label{probNot2}
Karena peluang memperoleh \resp{2} adalah $1/6$ atau $16.\bar{6}\%$, peluang tidak memperoleh \resp{2} adalah $100\% - 16.\bar{6}\%=83.\bar{3}\%$ atau $5/6$.

Sebagai alternatif, kita dapat menyadari bahwa tidak memperoleh \resp{2} sama dengan memperoleh \resp{1}, \resp{3}, \resp{4}, \resp{5}, atau \resp{6}. Kelima hasil itu mencakup lima dari enam hasil yang sama-sama berpeluang, sehingga probabilitasnya $5/6$.
\end{nexample}
\end{examplewrap}

\begin{examplewrap}
\begin{nexample}{Misalkan kita melempar dua dadu. Jika pada $1/6$ bagian dari seluruh lemparan dadu pertama menunjukkan \resp{1}, dan pada $1/6$ bagian dari waktu tersebut dadu kedua juga menunjukkan \resp{1}, berapa peluang memperoleh dua \resp{1}?}\label{probOf2Ones}
Jika dadu pertama menunjukkan \resp{1} pada $16.\bar{6}$\% lemparan dan pada $1/6$ bagian dari \emph{lemparan tersebut} dadu kedua juga menunjukkan \resp{1}, maka peluang kedua dadu menunjukkan \resp{1} adalah $(1/6)\times (1/6)$ atau $1/36$.
\end{nexample}
\end{examplewrap}


\D{\newpage}

\subsection{Probabilitas}

\index{random process|(}

Kita menggunakan probabilitas untuk membangun alat guna menggambarkan dan memahami keacakan yang tampak. Kita sering merumuskan probabilitas melalui \term{proses acak} yang menghasilkan suatu \term{hasil}.
\begin{center}
\begin{tabular}{lll}
Lempar dadu &$\rightarrow$ & \resp{1}, \resp{2}, \resp{3}, \resp{4}, \resp{5}, atau \resp{6} \\
Lempar koin &$\rightarrow$ & \resp{H} atau \resp{T} \\
\end{tabular}
\end{center}
Melempar dadu atau koin merupakan proses yang tampak acak, dan masing-masing menghasilkan suatu hasil.

\begin{onebox}{Probabilitas}
\term{Probabilitas} suatu hasil adalah proporsi kemunculan hasil tersebut jika kita mengamati proses acak itu sebanyak tak berhingga kali.
\end{onebox}

Probabilitas didefinisikan sebagai proporsi dan selalu bernilai antara 0~dan~1, termasuk kedua batasnya. Probabilitas juga dapat ditampilkan sebagai persentase antara 0\% dan 100\%.

Probabilitas dapat diilustrasikan dengan melempar dadu berkali-kali. Misalkan $\hat{p}_n$ adalah proporsi hasil yang berupa \resp{1} setelah $n$ lemparan pertama. Ketika jumlah lemparan bertambah, $\hat{p}_n$ akan konvergen ke probabilitas memperoleh \resp{1}, yaitu $p = 1/6$. Gambar~\ref{dieProp} menunjukkan konvergensi ini untuk 100.000 lemparan dadu. Kecenderungan $\hat{p}_n$ untuk stabil di sekitar $p$ dijelaskan oleh \term{Hukum Bilangan Besar}.

\begin{figure}[h]
\centering
\Figure[Diagram garis ditampilkan. Sumbu horizontal adalah "n (jumlah lemparan)", dengan nilai yang meningkat secara eksponensial dari 1 ke 10, 100, 1.000, 10.000, lalu 100.000. Sumbu vertikal adalah "p-hat sub n" dengan rentang 0,0 hingga sekitar 0,35. Garis putus-putus horizontal juga ditampilkan pada nilai satu per enam. Garis yang menunjukkan fraksi lemparan bernilai 1 mulai dari 0 pada lemparan pertama dan tetap di sana sampai sekitar lemparan ke-4, lalu melonjak ke 0,25 dan berosilasi, kemudian mencapai sekitar 0,35 pada 10 lemparan sebelum turun mendekati satu per enam. Sampai 100 lemparan, garis berosilasi antara 0,13 dan 0,22; dengan semakin banyak lemparan, garis makin stabil di sekitar satu per enam, menyimpang tidak lebih dari sekitar 0,03 sampai 1.000 lemparan. Sampai sekitar 5.000 lemparan penyimpangannya tidak lebih dari sekitar 0,015, dan setelah itu hampir tidak dapat dibedakan dari satu per enam selama lebih dari 5.000 lemparan.]{0.85}{dieProp}
\caption{Fraksi lemparan dadu yang menghasilkan 1 pada setiap tahap simulasi. Seiring bertambahnya jumlah lemparan, proporsinya cenderung mendekati probabilitas $1/6 \approx 0.167$.}
\label{dieProp}
\end{figure}

\begin{onebox}{Hukum Bilangan Besar}
Semakin banyak pengamatan dikumpulkan, proporsi $\hat{p}_n$ kemunculan suatu hasil tertentu akan konvergen ke probabilitas $p$ hasil tersebut.
\end{onebox}

Sesekali proporsi akan menyimpang dari probabilitas dan tampak menentang Hukum Bilangan Besar, seperti yang sering terjadi pada $\hat{p}_n$ di Gambar~\ref{dieProp}. Namun, penyimpangan itu mengecil ketika jumlah lemparan bertambah.

Di atas kita menulis $p$ sebagai probabilitas memperoleh \resp{1}. Probabilitas ini juga dapat ditulis sebagai
\begin{align*}
P(\text{memperoleh \resp{1}})
\end{align*}
Seiring terbiasanya kita dengan notasi ini, kita dapat menyingkatnya lebih lanjut. Misalnya, jika prosesnya jelas berupa “melempar dadu”, kita dapat menyingkat $P($memperoleh \resp{1}$)$ menjadi~$P($\resp{1}$)$.

\begin{exercisewrap}
\begin{nexercise} \label{randomProcessExercise}
Contoh proses acak adalah melempar dadu dan melempar koin. (a) Pikirkan proses acak lain. (b) Uraikan semua hasil yang mungkin dari proses itu. Sebagai contoh, melempar dadu adalah proses acak dengan hasil yang mungkin \mbox{\resp{1}, \resp{2}, ..., \resp{6}}.\footnotemark
\end{nexercise}
\end{exercisewrap}
\footnotetext{Berikut empat contoh. (i) Apakah seseorang akan sakit bulan depan atau tidak merupakan proses yang tampak acak, dengan hasil \resp{sakit} dan \resp{tidak}. (ii) Kita dapat \emph{membangkitkan} proses acak dengan memilih seseorang secara acak lalu mengukur tinggi badannya. Hasil proses ini berupa bilangan positif. (iii) Apakah pasar saham naik atau turun minggu depan merupakan proses yang tampak acak, dengan hasil yang mungkin \resp{naik}, \resp{turun}, dan \resp{tidak\us{}berubah}. Sebagai alternatif, perubahan persentase pasar saham dapat digunakan sebagai hasil numerik. (iv) Apakah teman sekamar Anda mencuci piring malam ini juga mungkin tampak seperti proses acak, dengan hasil \resp{mencuci\us{}piring} dan \resp{meninggalkan\us{}piring}.}

Hal-hal yang kita anggap sebagai proses acak belum tentu benar-benar acak; mungkin saja proses itu terlalu sulit dipahami secara tepat. Contoh keempat dalam jawaban catatan kaki Latihan Terbimbing~\ref{randomProcessExercise} mengisyaratkan bahwa perilaku teman sekamar adalah proses acak. Namun, meskipun perilaku teman sekamar tidak benar-benar acak, memodelkannya sebagai proses acak tetap dapat berguna.

%\begin{tipBox}{\tipBoxTitle{Modeling a process as random}
%It can be helpful to model a process as random even if it is not truly random.}
%\end{tipBox}

\index{random process|)}

\subsection{Hasil saling lepas atau saling eksklusif}

\index{disjoint|(}
\index{mutually exclusive|(}

Dua hasil disebut \term{saling lepas} atau \term{saling eksklusif} jika keduanya tidak mungkin terjadi bersamaan. Misalnya, jika kita melempar dadu, hasil \resp{1} dan \resp{2} saling lepas karena keduanya tidak mungkin muncul sekaligus. Sebaliknya, hasil \resp{1} dan “memperoleh bilangan ganjil” tidak saling lepas karena keduanya terjadi jika hasil lemparan adalah \resp{1}. Istilah \emph{saling lepas} dan \emph{saling eksklusif} ekuivalen dan dapat dipertukarkan.

Menghitung probabilitas hasil yang saling lepas cukup mudah. Saat melempar dadu, hasil \resp{1} dan \resp{2} saling lepas, sehingga probabilitas salah satu hasil itu terjadi dihitung dengan menjumlahkan probabilitas masing-masing:
\begin{align*}
P(\text{\resp{1} atau \resp{2}})
  = P(\text{\resp{1}})+P(\text{\resp{2}})
  = 1/6 + 1/6
  = 1/3
\end{align*}
Bagaimana dengan probabilitas memperoleh \resp{1}, \resp{2}, \resp{3}, \resp{4}, \resp{5}, atau \resp{6}? Sekali lagi, semua hasil tersebut saling lepas, sehingga probabilitasnya dijumlahkan:
\begin{align*}
&P(\text{\resp{1} atau \resp{2} atau \resp{3} atau \resp{4}
    atau \resp{5} atau \resp{6}}) \\
  &\quad = P(\text{\resp{1}})+P(\text{\resp{2}})
      + P(\text{\resp{3}})+P(\text{\resp{4}})
      + P(\text{\resp{5}})+P(\text{\resp{6}}) \\
  &\quad = 1/6 + 1/6 + 1/6 + 1/6 + 1/6 + 1/6
  = 1
\end{align*}
\term{Aturan Penjumlahan} menjamin ketepatan pendekatan ini ketika hasil-hasilnya saling lepas.

\begin{onebox}{Aturan Penjumlahan untuk hasil yang saling lepas}
  Jika $A_1$ dan $A_2$ menyatakan dua hasil yang saling lepas,
  probabilitas salah satunya terjadi diberikan oleh
  \begin{align*}
  P(A_1\text{ atau } A_2) = P(A_1) + P(A_2)
  \end{align*}
  Jika terdapat banyak hasil yang saling lepas, $A_1$, ..., $A_k$,
  probabilitas salah satu hasil tersebut terjadi adalah
  \begin{align*}
  P(A_1) + P(A_2) + \cdots + P(A_k)
  \end{align*}
\end{onebox}

\D{\newpage}

\begin{exercisewrap}
\begin{nexercise}
Kita tertarik pada probabilitas memperoleh \resp{1}, \resp{4}, atau \resp{5}. (a) Jelaskan mengapa hasil \resp{1}, \resp{4}, dan \resp{5} saling lepas. (b) Terapkan Aturan Penjumlahan untuk hasil yang saling lepas guna menentukan $P($\resp{1} atau \resp{4} atau \resp{5}$)$.\footnotemark
\end{nexercise}
\end{exercisewrap}
\footnotetext{(a) Proses acaknya adalah lemparan dadu, dan paling banyak satu dari hasil tersebut yang dapat muncul. Jadi, hasil-hasil itu saling lepas. (b)~$P($\resp{1} atau \resp{4} atau \resp{5}$) = P($\resp{1}$)+P($\resp{4}$)+P($\resp{5}$) = \frac{1}{6} + \frac{1}{6} + \frac{1}{6} = \frac{3}{6} = \frac{1}{2}$}

\index{data!loans|(}
\begin{exercisewrap}
\begin{nexercise}
Dalam himpunan data \data{loans} pada Bab~\ref{ch_summarizing_data},
variabel \var{homeownership} menjelaskan apakah peminjam
menyewa, memiliki hipotek, atau memiliki rumah sendiri.
Dari 10.000 peminjam, 3858 menyewa, 4789 memiliki hipotek,
dan 1353 memiliki rumah sendiri.\footnotemark{}
\begin{enumerate}[(a)]
\setlength{\itemsep}{0mm}
\item
    Apakah hasil \resp{rent}, \resp{mortgage},
    dan \resp{own} saling lepas?
\item
    Tentukan proporsi pinjaman dengan nilai \resp{mortgage}
    dan \resp{own} secara terpisah.
\item
    Gunakan Aturan Penjumlahan untuk hasil yang saling lepas
    guna menghitung probabilitas bahwa pinjaman yang dipilih
    secara acak dari himpunan data tersebut adalah milik seseorang
    yang memiliki hipotek atau rumah sendiri.
\end{enumerate}
\end{nexercise}
\end{exercisewrap}
\footnotetext{(a)~Ya. Setiap pinjaman dikategorikan ke tepat satu
level \var{homeownership}.
(b)~Hipotek: $\frac{4789}{10000} = 0.479$.
Milik sendiri: $\frac{1353}{10000} = 0.135$.
(c)~$P($\resp{mortgage} atau \resp{own}$) = P($\resp{mortgage}$) + P($\resp{own}$) = 0.479 + 0.135 = 0.614$.}
\index{data!loans|)}

\index{event|(}

Ilmuwan data jarang bekerja dengan hasil individual; mereka lebih sering mempertimbangkan \indexthis{\emph{himpunan}}{sets} atau \indexthis{\emph{kumpulan}}{collections} hasil. Misalkan $A$ menyatakan peristiwa ketika lemparan dadu menghasilkan \resp{1} atau \resp{2}, dan $B$ menyatakan peristiwa ketika lemparan menghasilkan \resp{4} atau \resp{6}. Kita menulis $A$ sebagai himpunan hasil $\{$\resp{1},~\resp{2}$\}$ dan $B=\{$\resp{4}, \resp{6}$\}$. Himpunan-himpunan ini lazim disebut \termsub{peristiwa}{event}. Karena $A$ dan $B$ tidak memiliki elemen yang sama, keduanya adalah peristiwa yang saling lepas. $A$ dan $B$ ditunjukkan pada Gambar~\ref{disjointSets}.

\begin{figure}[hhh]
  \centering
  \Figure[Enam bilangan ditampilkan berurutan: 1, 2, 3, 4, 5, dan 6. Bilangan 1 dan 2 dilingkari dan diberi label huruf “A”; bilangan 2 dan 3 dilingkari dan diberi label huruf “B”; bilangan 4 dan 6 dilingkari dengan label huruf “C”. (Lingkaran terakhir bukan lingkaran sebenarnya, melainkan kurungan yang tidak mencakup bilangan 5.)]{0.45}{disjointSets}
  \caption{Tiga peristiwa, $A$, $B$, dan $D$, terdiri atas
      hasil dari lemparan dadu.
      $A$ dan $B$ saling lepas karena tidak memiliki
      hasil yang sama.}
  \label{disjointSets}
\end{figure}

Aturan Penjumlahan berlaku untuk hasil maupun peristiwa yang saling lepas. Probabilitas salah satu peristiwa saling lepas $A$ atau $B$ terjadi adalah jumlah probabilitas masing-masing:
\begin{align*}
P(A\text{ atau }B) = P(A) + P(B) = 1/3 + 1/3 = 2/3
\end{align*}

\begin{exercisewrap}
\begin{nexercise}
(a) Verifikasi bahwa probabilitas peristiwa $A$, $P(A)$,
adalah $1/3$ menggunakan Aturan Penjumlahan.
(b)~Lakukan hal yang sama untuk peristiwa $B$.\footnotemark
\end{nexercise}
\end{exercisewrap}
\footnotetext{(a)~$P(A) = P($\resp{1} atau \resp{2}$)
    = P($\resp{1}$) + P($\resp{2}$)
    = \frac{1}{6} + \frac{1}{6}
    = \frac{2}{6}
    = \frac{1}{3}$.
    (b)~Dengan cara serupa, $P(B) = 1/3$.}

\begin{exercisewrap}
\begin{nexercise} \label{exerExaminingDisjointSetsABD}
(a) Dengan menggunakan Gambar~\ref{disjointSets} sebagai acuan, hasil apa yang direpresentasikan oleh peristiwa $D$? (b) Apakah peristiwa $B$ dan $D$ saling lepas? (c) Apakah peristiwa $A$ dan $D$ saling lepas?\footnotemark
\end{nexercise}
\end{exercisewrap}
\footnotetext{(a)~Hasil \resp{2} dan \resp{3}. (b)~Ya, peristiwa $B$ dan $D$ saling lepas karena tidak memiliki hasil yang sama. (c)~Peristiwa $A$ dan $D$ memiliki hasil yang sama, yaitu \resp{2}, sehingga tidak saling lepas.}

\begin{exercisewrap}
\begin{nexercise}
Dalam Latihan Terbimbing~\ref{exerExaminingDisjointSetsABD}, Anda mengonfirmasi bahwa $B$ dan $D$ pada Gambar~\ref{disjointSets} saling lepas. Hitung probabilitas peristiwa $B$ atau $D$ terjadi.\footnotemark
\end{nexercise}
\end{exercisewrap}
\footnotetext{Karena $B$ dan $D$ adalah peristiwa yang saling lepas, gunakan Aturan Penjumlahan: $P(B$ atau $D) = P(B) + P(D) = \frac{1}{3} + \frac{1}{3} = \frac{2}{3}$.}

\index{event|)}
\index{disjoint|)}
\index{mutually exclusive|)}


\D{\newpage}

\subsection{Probabilitas ketika peristiwa tidak saling lepas}

Mari kita pertimbangkan perhitungan untuk dua peristiwa yang tidak saling lepas dalam konteks \indexthis{dek standar 52 kartu}{deck of cards}, yang ditampilkan pada Gambar~\ref{deckOfCards}. Jika Anda belum mengenal kartu dalam dek standar, lihat catatan kaki.\footnote{Ke-52 kartu dibagi menjadi empat \term{jenis}: $\clubsuit$ (keriting), {\color{redcards}$\diamondsuit$} (wajik), {\color{redcards}$\heartsuit$} (hati), dan $\spadesuit$ (sekop). Setiap jenis memiliki 13 kartu berlabel: \resp{2}, \resp{3}, ..., \resp{10}, \resp{J} (jack), \resp{Q} (queen), \resp{K} (king), dan \resp{A} (ace). Jadi, setiap kartu merupakan kombinasi unik antara jenis dan label, misalnya {\color{redcards}\resp{4$\heartsuit$}} dan \resp{J$\clubsuit$}. Ke-12 kartu berupa jack, queen, dan king disebut \termsub{\resp{kartu bergambar}}{face card}. Kartu berjenis {\color{redcards}$\diamondsuit$} atau {\color{redcards}$\heartsuit$} biasanya berwarna {\color{redcards}merah}, sedangkan dua jenis lainnya biasanya berwarna hitam.}

\begin{figure}[h]
\centering
\begin{tabular}{lll lll lll lll l}
\resp{2$\clubsuit$} & \resp{3$\clubsuit$} & \resp{4$\clubsuit$} & \resp{5$\clubsuit$} & \resp{6$\clubsuit$} & \resp{7$\clubsuit$} & \resp{8$\clubsuit$} & \resp{9$\clubsuit$} & \resp{10$\clubsuit$} & \resp{J$\clubsuit$} & \resp{Q$\clubsuit$} & \resp{K$\clubsuit$} & \resp{A$\clubsuit$}  \\
\color{redcards} \resp{2$\diamondsuit$} & \color{redcards}\resp{3$\diamondsuit$} & \color{redcards}\resp{4$\diamondsuit$} & \color{redcards}\resp{5$\diamondsuit$} & \color{redcards}\resp{6$\diamondsuit$} & \color{redcards}\resp{7$\diamondsuit$} & \color{redcards}\resp{8$\diamondsuit$} & \color{redcards}\resp{9$\diamondsuit$} & \color{redcards}\resp{10$\diamondsuit$} & \color{redcards}\resp{J$\diamondsuit$} & \color{redcards}\resp{Q$\diamondsuit$} & \color{redcards}\resp{K$\diamondsuit$} & \color{redcards}\resp{A$\diamondsuit$} \\
\color{redcards}\resp{2$\heartsuit$} & \color{redcards}\resp{3$\heartsuit$} & \color{redcards}\resp{4$\heartsuit$} & \color{redcards}\resp{5$\heartsuit$} & \color{redcards}\resp{6$\heartsuit$} & \color{redcards}\resp{7$\heartsuit$} & \color{redcards}\resp{8$\heartsuit$} & \color{redcards}\resp{9$\heartsuit$} & \color{redcards}\resp{10$\heartsuit$} & \color{redcards}\resp{J$\heartsuit$} & \color{redcards}\resp{Q$\heartsuit$} & \color{redcards}\resp{K$\heartsuit$} & \color{redcards}\resp{A$\heartsuit$} \\
\resp{2$\spadesuit$} & \resp{3$\spadesuit$} & \resp{4$\spadesuit$} & \resp{5$\spadesuit$} & \resp{6$\spadesuit$} & \resp{7$\spadesuit$} & \resp{8$\spadesuit$} & \resp{9$\spadesuit$} & \resp{10$\spadesuit$} & \resp{J$\spadesuit$} & \resp{Q$\spadesuit$} & \resp{K$\spadesuit$} & \resp{A$\spadesuit$}
\end{tabular}
\caption{Representasi 52 kartu unik dalam sebuah dek.}
\label{deckOfCards}
\end{figure}

\begin{exercisewrap}
\begin{nexercise}
(a) Berapa probabilitas kartu yang dipilih secara acak merupakan wajik? (b)~Berapa probabilitas kartu yang dipilih secara acak merupakan kartu bergambar?\footnotemark
\end{nexercise}
\end{exercisewrap}
\footnotetext{(a) Ada 52 kartu dan 13 kartu wajik. Jika kartu dikocok dengan baik, setiap kartu memiliki peluang yang sama untuk terambil, sehingga probabilitas kartu yang dipilih secara acak merupakan wajik adalah $P({\color{redcards}\diamondsuit}) = \frac{13}{52} = 0.250$. (b)~Demikian pula, ada 12 kartu bergambar, sehingga $P($kartu bergambar$) = \frac{12}{52} = \frac{3}{13} = 0.231$.}

\term{Diagram Venn} berguna ketika hasil dapat dikategorikan sebagai “masuk” atau “tidak masuk” untuk dua atau tiga variabel, atribut, atau proses acak. Diagram Venn pada Gambar~\ref{cardsDiamondFaceVenn} menggunakan satu lingkaran untuk mewakili wajik dan lingkaran lain untuk mewakili kartu bergambar. Jika sebuah kartu merupakan wajik sekaligus kartu bergambar, kartu itu berada di irisan kedua lingkaran. Jika kartu itu wajik tetapi bukan kartu bergambar, kartu itu berada di bagian lingkaran kiri yang tidak termasuk lingkaran kanan (dan seterusnya). Jumlah kartu wajik diberikan oleh seluruh kartu di dalam lingkaran wajik: $10+3=13$. Probabilitasnya juga ditampilkan (misalnya, $10/52 = 0.1923$).

\begin{figure}[h]
\centering
\Figure[Diagram Venn ditampilkan. Satu lingkaran berlabel “Wajik” dengan proporsi total 0,25 dan lingkaran kedua berlabel “Kartu bergambar” dengan proporsi total 0,2308. Kedua lingkaran beririsan dan berbagi 3 kartu, dengan proporsi 0,0577 dari satu dek. Bagian lingkaran wajik yang tidak beririsan berlabel “10” untuk 10 kartu dan proporsi 0,1923. Bagian lingkaran kartu bergambar yang tidak beririsan berlabel “9” untuk 9 kartu dan proporsi 0,2308. Gambar juga mencatat bahwa terdapat 30 kartu yang bukan wajik maupun kartu bergambar.]{0.65}{cardsDiamondFaceVenn}
\caption{Diagram Venn untuk wajik dan kartu bergambar.}
\label{cardsDiamondFaceVenn}
\end{figure}

Misalkan $A$ menyatakan peristiwa bahwa kartu yang dipilih secara acak adalah wajik dan $B$ menyatakan peristiwa bahwa kartu tersebut merupakan kartu bergambar. Bagaimana menghitung $P(A$ atau $B)$? Peristiwa $A$ dan $B$ tidak saling lepas--kartu {\color{redcards}$J\diamondsuit$}, {\color{redcards}$Q\diamondsuit$}, dan {\color{redcards}$K\diamondsuit$} termasuk dalam kedua kategori--sehingga kita tidak dapat menggunakan Aturan Penjumlahan untuk peristiwa yang saling lepas. Sebagai gantinya, kita menggunakan diagram Venn. Kita mulai dengan menjumlahkan probabilitas kedua peristiwa:
\begin{align*}
P(A) + P(B)
  = P({\color{redcards}\diamondsuit}) + P(\text{kartu bergambar})
  = 13/52 + 12/52
\end{align*}

\D{\newpage}

\noindent%
Namun, tiga kartu yang termasuk dalam kedua peristiwa telah dihitung dua kali, sekali pada setiap probabilitas. Kita harus memperbaiki penghitungan ganda ini:
\begin{align*}
P(A\text{ atau } B)
  &= P({\color{redcards}\diamondsuit}\text{ atau kartu bergambar}) \\
  &= P({\color{redcards}\diamondsuit}) + P(\text{kartu bergambar})
      - P({\color{redcards}\diamondsuit}\text{ dan kartu bergambar}) \\
  &= 13/52 + 12/52 - 3/52 \\
  &= 22/52 = 11/26
\end{align*}
Persamaan ini merupakan contoh \term{Aturan Penjumlahan Umum}.

\begin{onebox}{Aturan Penjumlahan Umum}
  Jika $A$ dan $B$ adalah sembarang dua peristiwa, baik saling lepas maupun tidak,
  probabilitas bahwa setidaknya salah satunya terjadi adalah
  \begin{align*}
  P(A\text{ atau }B) = P(A) + P(B) - P(A\text{ dan }B)
  \end{align*}
  dengan $P(A$ dan $B)$ adalah probabilitas kedua peristiwa terjadi.
\end{onebox}

\begin{tipBox}{\tipBoxTitle{“atau” bersifat inklusif}
Ketika kita menulis “atau” dalam statistika, yang dimaksud adalah “dan/atau” kecuali dinyatakan secara eksplisit. Jadi, $A$ atau $B$ terjadi berarti $A$, $B$, atau keduanya, $A$ dan $B$, terjadi.}
\end{tipBox}

\begin{exercisewrap}
\begin{nexercise}
(a) Jika $A$ dan $B$ saling lepas, jelaskan mengapa hal ini berarti $P(A$ dan $B) = 0$. (b) Dengan menggunakan bagian (a), verifikasi bahwa Aturan Penjumlahan Umum menyederhana menjadi Aturan Penjumlahan untuk peristiwa saling lepas jika $A$ dan $B$ saling lepas.\footnotemark
\end{nexercise}
\end{exercisewrap}
\footnotetext{(a) Jika $A$ dan $B$ saling lepas, $A$ dan $B$ tidak mungkin terjadi secara bersamaan. (b) Jika $A$ dan $B$ saling lepas, suku terakhir $P(A\text{ dan }B)$ dalam rumus Aturan Penjumlahan Umum bernilai 0 (lihat bagian (a)), sehingga tersisa Aturan Penjumlahan untuk peristiwa saling lepas.}

\index{data!loans|(}

\begin{exercisewrap}
\begin{nexercise}\label{emailSpamNumberVennExer}
% library(openintro); d <- loans_full_schema; table(d[,c("application_type", "homeownership")]); table(d[,c("application_type")]); table(d[,c("homeownership")])
Dalam himpunan data \data{loans} yang memuat 10.000 pinjaman,
1495 pinjaman berasal dari permohonan bersama
(misalnya, pasangan mengajukan bersama),
4789 pemohon memiliki hipotek,
dan 950 memiliki kedua karakteristik tersebut.
Buat diagram Venn untuk situasi ini.\footnotemark
\end{nexercise}
\end{exercisewrap}
\footnotetext{%
  \begin{minipage}[t]{0.65\textwidth}
  Baik jumlah maupun {\color{oiB}probabilitas} yang bersesuaian
  (misalnya, $3839/10000 = 0.384$) ditampilkan.
  Perhatikan bahwa jumlah pinjaman pada lingkaran kiri
  adalah $3839 + 950 = 4789$, sedangkan jumlah
  pada lingkaran kanan adalah $950 + 545 = 1495$.
  \end{minipage}\ %
  \begin{minipage}[c]{0.3\textwidth}
  \hfill\Figure[Diagram Venn dengan dua lingkaran ditampilkan. Lingkaran pertama berlabel “pemohon memiliki hipotek” dan lingkaran kedua berlabel “permohonan bersama”; kedua lingkaran sebagian beririsan. Pada lingkaran pertama, bagian yang tidak beririsan menunjukkan jumlah 3839 dan proporsi 0,384. Bagian lingkaran permohonan bersama yang tidak beririsan berlabel 545 dengan proporsi 0,055. Bagian irisan berlabel jumlah 950 dan proporsi 0,095. Di luar kedua lingkaran dicatat bahwa “pinjaman lain” berjumlah 10.000 dikurangi 3.839, 950, dan 545, yaitu 4666 dengan proporsi 0,467.]{}{loans_app_type_home_venn} \vspace{-13mm}
  \end{minipage}}

\begin{exercisewrap}
\begin{nexercise}
(a)~Gunakan diagram Venn dari Latihan Terbimbing~\ref{emailSpamNumberVennExer} untuk menentukan
probabilitas bahwa pinjaman yang dipilih secara acak dari himpunan data \data{loans}
berasal dari permohonan bersama dengan pasangan yang memiliki
hipotek.
(b)~Berapa probabilitas pinjaman memiliki salah satu
atribut tersebut?\footnotemark
\end{nexercise}
\end{exercisewrap}
\footnotetext{%
  (a)~Jawabannya direpresentasikan oleh irisan
  kedua lingkaran: 0.095.
  (b)~Ini adalah jumlah tiga probabilitas yang saling lepas yang ditunjukkan
  di dalam lingkaran: $0.384 + 0.095 + 0.055 = 0.534$
  (berbeda 0.001 karena pembulatan).}

\index{data!loans|)}


\D{\newpage}

\subsection{Distribusi probabilitas}

\termsub{Distribusi probabilitas}{probability!distribution} adalah tabel semua hasil yang saling lepas beserta probabilitasnya. Gambar~\ref{diceProb} menunjukkan distribusi probabilitas jumlah dua dadu.

\begin{figure}[h] \small
\centering
\begin{tabular}{l ccc ccc ccc cc}
  \hline
  \ \vspace{-3mm} \\
Jumlah dadu\vspace{0.3mm} & 2 & 3 & 4 & 5 & 6 & 7 & 8 & 9 & 10 & 11 & 12  \\
Probabilitas & $\frac{1}{36}$ & $\frac{2}{36}$ & $\frac{3}{36}$ & $\frac{4}{36}$ & $\frac{5}{36}$ & $\frac{6}{36}$ & $\frac{5}{36}$ & $\frac{4}{36}$ & $\frac{3}{36}$ & $\frac{2}{36}$ & $\frac{1}{36}$\vspace{1mm} \\
   \hline
\end{tabular}
\caption{Distribusi probabilitas jumlah dua dadu.}
\label{diceProb}
\end{figure}

\begin{onebox}{Aturan distribusi probabilitas}
Distribusi probabilitas adalah daftar hasil yang mungkin beserta probabilitasnya, yang memenuhi tiga aturan berikut: \vspace{-2mm}
\begin{enumerate}
\setlength{\itemsep}{0mm}
\item Hasil yang tercantum harus saling lepas.
\item Setiap probabilitas harus berada antara 0 dan 1.
\item Jumlah seluruh probabilitas harus 1. \vspace{1mm}
\end{enumerate}
\end{onebox}

\begin{exercisewrap}
\begin{nexercise}\label{usHouseholdIncomeDistsExercise}
Gambar~\ref{usHouseholdIncomeDists} menyajikan tiga distribusi pendapatan rumah tangga di Amerika Serikat. Hanya satu yang benar. Yang mana? Apa kesalahan pada dua lainnya?\footnotemark
\end{nexercise}
\end{exercisewrap}
\footnotetext{Probabilitas pada (a) tidak berjumlah~1.
    Probabilitas kedua pada (b) bernilai negatif.
    Dengan demikian tersisa (c), yang memang memenuhi
    syarat sebuah distribusi.
    Salah satu dari ketiganya disebut sebagai
    distribusi pendapatan rumah tangga AS yang sebenarnya,
    sehingga jawabannya pasti (c).}

\begin{figure}[h]
\centering
\begin{tabular}{r | cc cc}
  \hline
Rentang pendapatan & \$0-25k & \$25k-50k & \$50k-100k & \$100k+ \\
  \hline
(a)\hspace{0.2mm}	 & 0.18 & 0.39 & 0.33 & 0.16 \\
(b)				 & 0.38 & -0.27 & 0.52 & 0.37 \\
(c)\hspace{0.2mm}	 & 0.28 & 0.27 & 0.29 & 0.16 \\
  \hline
\end{tabular}
\caption{Distribusi pendapatan rumah tangga AS yang diusulkan (Latihan Terbimbing~\ref{usHouseholdIncomeDistsExercise}).}
\label{usHouseholdIncomeDists}
\end{figure}

Bab~\ref{introductionToData} menekankan pentingnya membuat plot data untuk memperoleh ringkasan cepat. Distribusi probabilitas juga dapat diringkas dengan diagram batang. Misalnya, distribusi pendapatan rumah tangga AS ditampilkan sebagai diagram batang pada Gambar~\ref{usHouseholdIncomeDistBar}. %\footnote{Plot distribusi juga dapat dibuat ketika pendapatan tidak dibagi secara artifisial ke dalam empat kelompok. Distribusi \emph{kontinu} dibahas pada Bagian~\ref{contDist}.}
Distribusi probabilitas jumlah dua dadu ditampilkan pada Gambar~\ref{diceProb} dan diplot pada Gambar~\ref{diceSumDist}.

\begin{figure}[h]
  \centering
  \Figure[Diagram batang “Pendapatan Rumah Tangga AS” ditampilkan dengan empat kelompok pendapatan. Sumbu vertikal berlabel “Probabilitas”. Kelompok pertama adalah \$0 sampai \$25.000 dengan tinggi batang sekitar 0,28; kelompok kedua \$25.000 sampai \$50.000 dengan tinggi sekitar 0,27; kelompok ketiga \$50.000 sampai \$100.000 dengan tinggi sekitar 0,28; dan kelompok keempat di atas \$100.000 dengan tinggi sekitar 0,15.]{0.65}{usHouseholdIncomeDistBar}
  \caption{Distribusi probabilitas pendapatan rumah tangga AS.}
  \label{usHouseholdIncomeDistBar}
\end{figure}

\begin{figure}
  \centering
  \Figure[Diagram batang jumlah dua dadu ditampilkan; nilainya dapat berupa 2, 3, 4, 5, dan seterusnya sampai 12. Sumbu vertikal berlabel “Probabilitas”. Tinggi batang untuk 2 sekitar 0,025; 3 sebesar 0,055; 4 sebesar 0,09; 5 sebesar 0,115; 6 sebesar 0,14; 7 sebesar 0,165; 8 sebesar 0,14; 9 sebesar 0,115; 10 sebesar 0,09; 11 sebesar 0,055; dan 12 sebesar 0,025.]{0.67}{diceSumDist}
  \caption{Distribusi probabilitas jumlah dua dadu.}
  \label{diceSumDist}
\end{figure}

Pada diagram batang ini, tinggi batang merepresentasikan probabilitas hasil. Jika hasilnya numerik dan diskret, biasanya (secara visual) nyaman membuat diagram batang yang menyerupai histogram, seperti pada jumlah dua dadu. Contoh lain penempatan batang pada lokasi masing-masing ditunjukkan pada Gambar~\ref{bookCostDist} di halaman~\pageref{bookCostDist}.

\subsection{Komplemen suatu peristiwa}

Melempar dadu menghasilkan nilai dalam himpunan $\{$\resp{1}, \resp{2}, \resp{3}, \resp{4}, \resp{5}, \resp{6}$\}$. Himpunan semua hasil yang mungkin ini disebut \term{ruang sampel} ($S$)\index{S@$S$} untuk lemparan dadu. Kita sering menggunakan ruang sampel untuk memeriksa situasi ketika suatu peristiwa tidak terjadi.

Misalkan $D=\{$\resp{2}, \resp{3}$\}$ menyatakan peristiwa bahwa hasil lemparan dadu adalah \resp{2} atau \resp{3}. Kemudian \term{komplemen} $D$ menyatakan semua hasil dalam ruang sampel yang bukan anggota $D$, dilambangkan dengan $D^c = \{$\resp{1}, \resp{4}, \resp{5}, \resp{6}$\}$. Dengan kata lain, $D^c$ adalah himpunan semua hasil yang mungkin yang belum termasuk dalam $D$. Gambar~\ref{complementOfD} menunjukkan hubungan antara $D$, $D^c$, dan ruang sampel $S$.

\begin{figure}[hht]
  \centering
  \Figure[Bilangan 1, 2, 3, 4, 5, dan 6 ditampilkan berurutan. Bilangan 2 dan 3 dilingkari serta diberi label “D”. Bilangan 1, 4, 5, dan 6 dilingkari serta diberi label “D-pangkat-C” sebagai komplemen D. Semua bilangan kemudian dilingkari lagi dengan kurungan lebih besar berlabel “S” sebagai ruang sampel.]{0.55}{complementOfD}
  \caption{Peristiwa $D=\{$\resp{2}, \resp{3}$\}$ dan komplemennya,
      $D^c = \{$\resp{1}, \resp{4}, \resp{5}, \resp{6}$\}$.
      $S$~menyatakan ruang sampel, yaitu himpunan
      semua hasil yang mungkin.}
  \label{complementOfD}
\end{figure}

\begin{exercisewrap}
\begin{nexercise}
(a) Hitung $P(D^c) = P($memperoleh \resp{1}, \resp{4}, \resp{5}, atau \resp{6}$)$. (b) Berapakah $P(D) + P(D^c)$?\footnotemark
\end{nexercise}
\end{exercisewrap}
\footnotetext{(a)~Hasil-hasilnya saling lepas dan masing-masing memiliki probabilitas $1/6$, sehingga probabilitas totalnya $4/6=2/3$. (b)~Kita juga dapat melihat bahwa $P(D)=\frac{1}{6} + \frac{1}{6} = 1/3$. Karena $D$ dan $D^c$ saling lepas, $P(D) + P(D^c) = 1$.}

\begin{exercisewrap}
\begin{nexercise}
Peristiwa $A=\{$\resp{1}, \resp{2}$\}$ dan $B=\{$\resp{4}, \resp{6}$\}$ ditunjukkan pada Gambar~\ref{disjointSets} di halaman~\pageref{disjointSets}. (a) Tuliskan apa yang direpresentasikan oleh $A^c$ dan $B^c$. (b)~Hitung $P(A^c)$ dan $P(B^c)$. (c)~Hitung $P(A)+P(A^c)$ dan $P(B)+P(B^c)$.\footnotemark
\end{nexercise}
\end{exercisewrap}
\footnotetext{Jawaban singkat: (a)~$A^c=\{$\resp{3}, \resp{4}, \resp{5}, \resp{6}$\}$ dan $B^c=\{$\resp{1}, \resp{2}, \resp{3}, \resp{5}$\}$. (b)~Karena setiap hasil saling lepas, jumlahkan probabilitas hasil individual untuk memperoleh $P(A^c)=2/3$ dan $P(B^c)=2/3$. (c)~$A$~dan~$A^c$ saling lepas, demikian pula $B$~dan~$B^c$. Oleh karena itu, $P(A) + P(A^c) = 1$ dan $P(B) + P(B^c) = 1$.}

\D{\newpage}

Komplemen suatu peristiwa $A$ dibentuk agar memiliki dua sifat penting: (i) setiap hasil yang mungkin dan bukan anggota $A$ merupakan anggota $A^c$, dan (ii) $A$ dan $A^c$ saling lepas. Sifat (i) menyiratkan
\begin{align*}
P(A\text{ atau }A^c) = 1
\end{align*}
Artinya, jika hasil bukan anggota $A$, hasil itu harus direpresentasikan dalam $A^c$. Kita menggunakan Aturan Penjumlahan untuk peristiwa saling lepas guna menerapkan Sifat (ii):
\begin{align*}
P(A\text{ atau }A^c) = P(A) + P(A^c)
\end{align*}
Menggabungkan dua persamaan terakhir menghasilkan hubungan yang sangat berguna antara probabilitas suatu peristiwa dan komplemennya.

\begin{onebox}{Komplemen}
  Komplemen peristiwa $A$ dilambangkan $A^c$, dan $A^c$
  menyatakan semua hasil yang bukan anggota~$A$. $A$ dan $A^c$
  berhubungan secara matematis:
  \begin{align*}
  P(A) + P(A^c) = 1, \quad\text{i.e.}\quad P(A) = 1-P(A^c)
  \end{align*}
\end{onebox}

Dalam contoh sederhana, menghitung $A$ atau $A^c$ dapat dilakukan dalam beberapa langkah. Namun, menggunakan komplemen dapat menghemat banyak waktu ketika masalah menjadi lebih kompleks.

\begin{exercisewrap}
\begin{nexercise}
Misalkan $A$ menyatakan peristiwa ketika kita melempar dua dadu dan jumlahnya kurang dari \resp{12}. (a)~Apa yang direpresentasikan oleh peristiwa $A^c$? (b)~Tentukan $P(A^c)$ dari Gambar~\ref{diceProb} pada halaman~\pageref{diceProb}. (c)~Tentukan $P(A)$.\footnotemark
\end{nexercise}
\end{exercisewrap}
\footnotetext{(a)~Komplemen $A$: jumlahnya sama dengan \resp{12}. (b)~$P(A^c) = 1/36$. (c)~Gunakan probabilitas komplemen dari bagian (b), $P(A^c) = 1/36$, dan persamaan komplemen: $P($kurang dari \resp{12}$) = 1 - P($\resp{12}$) = 1 - 1/36 = 35/36$.}

\begin{exercisewrap}
\begin{nexercise}
Tentukan probabilitas berikut untuk lemparan dua dadu:\footnotemark
\begin{enumerate}[(a)]
\setlength{\itemsep}{0mm}
\item Jumlah dadu \emph{bukan} \resp{6}.
\item Jumlahnya paling sedikit \resp{4}.
    Artinya, tentukan probabilitas peristiwa
    $B = \{$\resp{4}, \resp{5}, ..., \resp{12}$\}$.
\item Jumlahnya paling banyak \resp{10}.
    Artinya, tentukan probabilitas peristiwa
    $D=\{$\resp{2}, \resp{3}, ..., \resp{10}$\}$.
\end{enumerate}
\end{nexercise}
\end{exercisewrap}
\footnotetext{(a)~Pertama, cari $P($\resp{6}$)=5/36$, lalu gunakan komplemen: $P($bukan \resp{6}$) = 1 - P($\resp{6}$) = 31/36$.

(b)~Pertama, cari komplemennya, yang membutuhkan jauh lebih sedikit usaha: $P($\resp{2} atau \resp{3}$)=1/36+2/36=1/12$. Kemudian hitung $P(B) = 1-P(B^c) = 1-1/12 = 11/12$.

(c)~Seperti sebelumnya, mencari komplemen adalah cara cerdas untuk menentukan $P(D)$. Pertama, cari $P(D^c) = P($\resp{11} atau \resp{12}$)=2/36 + 1/36=1/12$. Kemudian hitung $P(D) = 1 - P(D^c) = 11/12$.}


\subsection{Kemandirian}
\label{probabilityIndependence}

Sebagaimana variabel dan pengamatan dapat saling independen, proses acak juga dapat saling independen. Dua proses \term{independen} jika mengetahui hasil salah satunya tidak memberikan informasi berguna tentang hasil yang lain. Misalnya, melempar koin dan melempar dadu adalah dua proses independen--mengetahui koin menunjukkan kepala tidak membantu menentukan hasil lemparan dadu. Sebaliknya, harga saham biasanya bergerak naik atau turun bersama-sama, sehingga tidak independen.

Contoh~\ref{probOf2Ones} memberikan contoh dasar dua proses independen: melempar dua dadu. Kita ingin menentukan probabilitas keduanya menunjukkan \resp{1}. Misalkan salah satu dadu berwarna merah dan yang lain putih. Jika hasil dadu merah adalah \resp{1}, hal itu tidak memberikan informasi tentang hasil dadu putih. Kita telah menjumpai pertanyaan yang sama pada Contoh~\ref{probOf2Ones} (halaman~\pageref{probOf2Ones}), ketika probabilitas dihitung dengan penalaran berikut: dadu merah menunjukkan \resp{1} pada $1/6$ bagian waktu, dan pada $1/6$ bagian dari \emph{waktu tersebut} dadu putih juga menunjukkan \resp{1}. Hal ini diilustrasikan pada Gambar~\ref{indepForRollingTwo1s}. Karena lemparan-lemparan itu independen, probabilitas hasil yang bersesuaian dapat dikalikan untuk memperoleh jawaban akhir: $(1/6)\times(1/6)=1/36$. Penalaran ini dapat digeneralisasi ke banyak proses independen.

\begin{figure}[hht]
\centering
\Figure[Persegi panjang hitam menguraikan gambar dan berlabel “Semua lemparan”. Di dalamnya, pita vertikal selebar kira-kira seperenam diarsir dan diberi label “seperenam lemparan pertama bernilai 1”. Bagian horizontal yang mewakili kira-kira seperenam dari pita vertikal itu diarsir berbeda dan diberi label “pada seperenam waktu ketika lemparan pertama bernilai 1, lemparan kedua juga bernilai 1”.]{0.6}{indepForRollingTwo1s}
\caption{Pada $1/6$ bagian waktu, lemparan pertama bernilai \resp{1}. Kemudian pada $1/6$ bagian dari \emph{waktu tersebut}, lemparan kedua juga bernilai \resp{1}.}
\label{indepForRollingTwo1s}
\end{figure}

\begin{examplewrap}
\begin{nexample}{Bagaimana jika ada dadu biru yang juga independen dari dua dadu lainnya? Berapa probabilitas melempar ketiga dadu dan semuanya menunjukkan \resp{1}?}\label{threeDice}
Logika yang sama seperti pada Contoh~\ref{probOf2Ones} berlaku. Jika dadu putih dan merah sama-sama \resp{1} pada $1/36$ bagian waktu, maka pada $1/6$ bagian dari \emph{waktu tersebut} dadu biru juga akan \resp{1}, sehingga kita mengalikan:
{\begin{align*}
P(putih=\text{\small\resp{1} dan } merah=\text{\small\resp{1} dan } biru=\text{\small\resp{1}})
	&= P(putih=\text{\small\resp{1}})\times P(merah=\text{\small\resp{1}})\times P(biru=\text{\small\resp{1}}) \\
	&= (1/6)\times (1/6)\times (1/6)
	= 1/216
\end{align*}} \vspace{-7mm}
\end{nexample}
\end{examplewrap}

Contoh~\ref{threeDice} mengilustrasikan apa yang disebut Aturan Perkalian untuk proses independen.

\begin{onebox}{Aturan Perkalian untuk proses independen}
  \index{Multiplication Rule|textbf}%
  Jika $A$ dan $B$ menyatakan peristiwa dari dua proses
  yang berbeda dan independen, probabilitas bahwa $A$
  dan $B$ terjadi dapat dihitung sebagai hasil kali
  probabilitas masing-masing:
  \begin{align*}
  P(A \text{ dan }B) = P(A) \times  P(B)
  \end{align*}
  Demikian pula, jika terdapat $k$ peristiwa $A_1$, ..., $A_k$
  dari $k$ proses independen, probabilitas semuanya terjadi adalah
  \begin{align*}
  P(A_1) \times  P(A_2)\times  \cdots \times  P(A_k)
  \end{align*}\vspace{-6mm}
\end{onebox}

\begin{exercisewrap}
\begin{nexercise} \label{ex2Handedness}
Sekitar 9\% orang kidal. Misalkan 2 orang dipilih secara acak dari populasi AS. Karena ukuran sampel 2 sangat kecil dibandingkan populasi, wajar untuk menganggap kedua orang ini independen. (a)~Berapa probabilitas keduanya kidal? (b)~Berapa probabilitas keduanya tidak kidal?\footnotemark
\end{nexercise}
\end{exercisewrap}
\footnotetext{(a) Probabilitas orang pertama kidal adalah $0.09$, sama dengan orang kedua. Kita menerapkan Aturan Perkalian untuk proses independen untuk menentukan probabilitas keduanya kidal: $0.09\times 0.09 = 0.0081$.

(b) Wajar untuk menganggap proporsi orang ambidekster (dapat menggunakan tangan kanan dan kiri) mendekati 0, sehingga $P($tidak kidal$)=1-0.09=0.91$. Dengan penalaran yang sama seperti bagian~(a), probabilitas keduanya tidak kidal adalah $0.91\times 0.91 = 0.8281$.}

\begin{exercisewrap}
\begin{nexercise}
\label{ex5Handedness}%
Misalkan 5 orang dipilih secara acak.\footnotemark\vspace{-1.5mm}
\begin{enumerate}
\setlength{\itemsep}{0mm}
\item[(a)] Berapa probabilitas semuanya tidak kidal?
\item[(b)] Berapa probabilitas semuanya kidal?
\item[(c)] Berapa probabilitas tidak semua orang tidak kidal?
\end{enumerate}
\end{nexercise}
\end{exercisewrap}
\footnotetext{(a)~Singkatan \resp{RH} dan \resp{LH} digunakan masing-masing untuk tidak kidal dan kidal. Karena semuanya independen, kita menerapkan Aturan Perkalian untuk proses independen:
\begin{align*}
P(\text{kelima orang adalah \resp{RH}})
&= P(\text{pertama = \resp{RH}, kedua = \resp{RH}, ..., kelima = \resp{RH}}) \\
&= P(\text{pertama = \resp{RH}})\times P(\text{kedua = \resp{RH}})\times  \dots \times P(\text{kelima = \resp{RH}}) \\
&= 0.91\times 0.91\times 0.91\times 0.91\times 0.91 = 0.624
\end{align*}

(b)~Dengan penalaran yang sama seperti pada~(a), $0.09\times 0.09\times 0.09\times 0.09\times 0.09 = 0.0000059$

(c)~Gunakan komplemen, $P($kelima orang \resp{RH}$)$, untuk menjawab pertanyaan ini:
\begin{align*}
P(\text{not all \resp{RH}})
	= 1 - P(\text{all \resp{RH}})
	= 1 - 0.624 = 0.376
\end{align*}}

Misalkan variabel \var{handedness} dan
\var{sex} independen,
artinya mengetahui \var{sex} seseorang tidak memberikan informasi
berguna tentang \var{handedness}-nya, dan sebaliknya.
Kita dapat menghitung probabilitas orang yang dipilih secara acak
tidak kidal dan perempuan\footnote{Proporsi aktual
  populasi AS yang \resp{perempuan} sekitar 50\%,
  sehingga kita menggunakan 0,5 sebagai probabilitas memilih perempuan.
  Namun, probabilitas ini berbeda di negara lain.}
dengan Aturan Perkalian:
\begin{align*}
P(\text{tidak kidal dan perempuan})
    &= P(\text{tidak kidal}) \times  P(\text{perempuan}) \\
    &= 0.91 \times  0.50 = 0.455
\end{align*}


\begin{exercisewrap}
\begin{nexercise}
Tiga orang dipilih secara acak.\footnotemark \vspace{-1.5mm}
\begin{enumerate}
\setlength{\itemsep}{0mm}
\item[(a)] Berapa probabilitas orang pertama laki-laki dan tidak kidal?
\item[(b)] Berapa probabilitas dua orang pertama laki-laki dan tidak kidal?
\item[(c)] Berapa probabilitas orang ketiga perempuan dan kidal?
\item[(d)] Berapa probabilitas dua orang pertama laki-laki dan tidak kidal, sedangkan orang ketiga perempuan dan kidal?
\end{enumerate}
\end{nexercise}
\end{exercisewrap}
\footnotetext{Jawaban singkat diberikan. (a)~Dalam notasi probabilitas, ini dapat ditulis $P($orang yang dipilih secara acak adalah laki-laki dan tidak kidal$)=0.455$. (b)~0.207. (c)~0.045. (d)~0.0093.}

Kadang-kadang kita ingin tahu apakah suatu hasil memberikan informasi berguna tentang hasil lain. Pertanyaannya: apakah kemunculan kedua peristiwa itu independen? Dua peristiwa $A$ dan $B$ disebut independen jika memenuhi
$P(A \text{ dan }B) = P(A) \times  P(B)$.

\begin{examplewrap}
\begin{nexample}{Jika kita mengocok dek kartu lalu mengambil satu kartu, apakah peristiwa kartu itu hati independen dari peristiwa kartu itu ace?}
Probabilitas kartu itu hati adalah $1/4$ dan probabilitas kartu itu ace adalah $1/13$. Probabilitas kartu itu ace hati adalah $1/52$.
Kita memeriksa apakah $P(A \text{ dan }B) = P(A) \times  P(B)$
terpenuhi:
\begin{align*}
P({\color{redcards}\heartsuit})\times P(\text{ace}) = \frac{1}{4}\times \frac{1}{13} = \frac{1}{52} 
					= P({\color{redcards}\heartsuit}\text{ dan ace})
\end{align*}
Karena persamaan tersebut berlaku, peristiwa kartu itu hati dan peristiwa kartu itu ace merupakan peristiwa independen.
\end{nexample}
\end{examplewrap}

% B009 translated boundary ends here: Chapter 3 opener and Defining probability
% through Independence (source lines 1--684). The separate EOC-only input below
% remains an untranslated upstream witness and is intentionally outside this
% candidate; it has no public answer payload in this source file.
{\input{ch_probability/TeX/defining_probability.tex}}




%_________________
\section{Peluang bersyarat}
\label{conditionalProbabilitySection}

Hubungan antara dua variabel atau lebih dapat sangat kaya
dan penting untuk dipahami.
Sebagai contoh, perusahaan asuransi mobil akan mempertimbangkan
informasi tentang riwayat mengemudi seseorang untuk menilai
risiko bahwa orang tersebut akan menyebabkan kecelakaan.
Hubungan semacam ini merupakan ranah
peluang bersyarat.


\subsection{Mengeksplorasi peluang dengan tabel kontingensi}

\index{data!photo\_classify|(}

\newcommand{\fashN}{1822}
% Urutan: ML, lalu manusia
\newcommand{\fashYY}{197}
\newcommand{\fashYN}{22}
\newcommand{\fashYA}{219}
\newcommand{\fashNY}{112}
\newcommand{\fashNN}{1491}
\newcommand{\fashNA}{1603}
\newcommand{\fashAY}{309}
\newcommand{\fashAN}{1513}
\newcommand{\fashAA}{\fashN{}}
%\newcommand{\fashPYY}{}
%\newcommand{\fashPYN}{}
%\newcommand{\fashPNY}{}
%\newcommand{\fashPNN}{}
%\newcommand{\fashPYA}{0.12}
%\newcommand{\fashPNA}{0.88}
%\newcommand{\fashPAY}{}
%\newcommand{\fashPAN}{}
%\newcommand{\fashPYCY}{}
%\newcommand{\fashPYCN}{}
%\newcommand{\fashPNCY}{}
%\newcommand{\fashPNCN}{}
\newcommand{\fashCYPY}{0.96}
\newcommand{\fashCYPN}{0.04}
\newcommand{\fashCNPY}{0.07}
\newcommand{\fashCNPN}{0.93}

Himpunan data \data{photo\us{}classify} memuat hasil kerja
sebuah pengklasifikasi pada sampel \fashN{} foto dari situs berbagi foto.
Para ilmuwan data berupaya meningkatkan kemampuan pengklasifikasi
untuk menentukan apakah sebuah foto berkaitan dengan mode, dan 1.822 foto ini
digunakan untuk menguji pengklasifikasi tersebut.
Setiap foto memperoleh dua klasifikasi:
yang pertama disebut \var{mach\us{}learn} dan berisi
klasifikasi dari sistem pembelajaran
mesin~(ML)\index{machine learning (ML)}, yakni
\resp{pred\us{}fashion} atau \resp{pred\us{}not}.
Masing-masing dari \fashN{} foto ini juga telah diklasifikasikan dengan cermat
oleh sebuah tim manusia, yang kita jadikan acuan kebenaran;
variabel ini disebut \var{truth} dan bernilai
\resp{fashion} atau \resp{not}.
Gambar~\ref{contTableOfFashionPhotos} merangkum hasilnya.

\begin{figure}[ht]
\centering
\begin{tabular}{ll ccc rr}
&& \multicolumn{2}{c}{\var{truth}} & \hspace{1cm} &  \\
\cline{3-4}
&& \resp{fashion} & \resp{not} & Total  \\
\cline{2-5}
& \resp{pred\us{}fashion} &
    \fashYY{} & \fashYN{} & \fashYA{} \\
\raisebox{1.5ex}[0pt]{\var{mach\us{}learn}}
    & \resp{pred\us{}not} \hspace{0.5cm} &
    \fashNY{} & \fashNN{} & \fashNA{}   \\
\cline{2-5}
& Total & \fashAY{} & \fashAN{} & \fashN{} \\
\end{tabular}
\caption{Tabel kontingensi yang merangkum himpunan data
    \data{photo\us{}classify}.}
\label{contTableOfFashionPhotos}
\end{figure}
% library(openintro); table(photo_classify)

\begin{figure}[ht]
  \centering
  \Figure[Ditampilkan diagram Venn yang menggunakan kotak, bukan lingkaran, untuk dua kategori yang sebagian bertumpang tindih, yaitu "ML Memprediksi Mode" dan "Foto Mode". Bagian persegi panjang ML Memprediksi Mode yang tidak bertumpang tindih berlabel 0,01. Bagian persegi panjang Foto Mode yang tidak bertumpang tindih berlabel 0,06. Bagian yang bertumpang tindih berlabel 0,11. Di luar kedua persegi panjang terdapat label "Bukan Keduanya" dengan nilai 0,82.]{0.65}{photoClassifyVenn}
  \caption{Diagram Venn yang menggunakan kotak untuk himpunan data
      \data{photo\us{}classify}.}
  \label{photoClassifyVenn}
\end{figure}

\begin{examplewrap}
\begin{nexample}{Jika sebuah foto memang berkaitan dengan mode,
    berapa peluang pengklasifikasi ML mengidentifikasi dengan benar
    bahwa foto tersebut berkaitan dengan mode?}
  Kita dapat memperkirakan peluang ini menggunakan data.
  Dari \fashAY{} foto mode,
  algoritma ML mengklasifikasikan \fashYY{} foto dengan benar:
\begin{align*}
P(\text{\var{mach\us{}learn} adalah \resp{pred\us{}fashion}
    jika \var{truth} adalah \resp{fashion}})
  = \frac{\fashYY{}}{\fashAY{}}
  = 0.638
\end{align*}
\end{nexample}
\end{examplewrap}

\begin{examplewrap}
\begin{nexample}{Kita mengambil sampel sebuah foto dari himpunan data
    dan mengetahui bahwa algoritma ML memprediksi foto ini
    tidak berkaitan dengan mode.
    Berapa peluang prediksi tersebut keliru dan
    foto itu sebenarnya berkaitan dengan mode?}
  Jika pengklasifikasi ML menyatakan bahwa sebuah foto tidak berkaitan dengan mode,
  foto itu berasal dari baris kedua pada himpunan data.
  Dari~\fashNA{} foto tersebut, \fashNY{} sebenarnya
  berkaitan dengan mode:
\begin{align*}
P(\text{\var{truth} adalah \resp{fashion}
    jika \var{mach\us{}learn} adalah \resp{pred\us{}not}})
  = \frac{\fashNY{}}{\fashNA{}}
  = 0.070
\end{align*}
\end{nexample}
\end{examplewrap}

\subsection{Peluang marginal dan bersama}
\label{marginalAndJointProbabilities}

\index{marginal probability|(}
\index{joint probability|(}

Gambar~\ref{contTableOfFashionPhotos} memuat total baris dan
kolom untuk setiap variabel secara terpisah dalam himpunan data
\data{photo\us{}classify}.
Total-total ini menyatakan
\termsub{peluang marginal}{peluang marginal}
untuk sampel tersebut, yaitu peluang yang didasarkan pada
satu variabel tanpa memperhatikan variabel lain.
Sebagai contoh, peluang yang hanya didasarkan pada variabel
\var{mach\us{}learn} merupakan peluang marginal:
\begin{align*}
P(\text{\var{mach\us{}learn} adalah \resp{pred\us{}fashion}})
    = \frac{\fashYA{}}{\fashN{}}
    = 0.12
\end{align*}
Peluang hasil untuk dua variabel atau proses atau lebih
disebut
\termsub{peluang \mbox{bersama}}{peluang bersama}:
\begin{align*}
P(\text{\var{mach\us{}learn} adalah \resp{pred\us{}fashion}
    dan \var{truth} adalah \resp{fashion}})
  = \frac{\fashYY{}}{\fashN{}}
  = 0.11
\end{align*}
Dalam peluang bersama, kata ``dan'' lazim diganti dengan koma,
meskipun penggunaan kata ``dan'' maupun koma
sama-sama dapat diterima:
\begin{center}
$P(\text{\var{mach\us{}learn} adalah \resp{pred\us{}fashion},
    \var{truth} adalah \resp{fashion}})$ \\[2mm]
berarti sama dengan \\[2mm]
$P(\text{\var{mach\us{}learn} adalah \resp{pred\us{}fashion}
    dan \var{truth} adalah \resp{fashion}})$
\end{center}

\begin{onebox}{Peluang marginal dan bersama}
  Jika suatu peluang didasarkan pada satu variabel,
  peluang tersebut merupakan \emph{\hiddenterm{peluang marginal}}.
  Peluang hasil untuk dua variabel atau proses atau lebih
  disebut \emph{\hiddenterm{peluang bersama}}.
\end{onebox}

Kita menggunakan \term{proporsi tabel} untuk merangkum peluang bersama
pada sampel \data{photo\us{}classify}.
Proporsi ini dihitung dengan membagi setiap frekuensi dalam
Gambar~\ref{contTableOfFashionPhotos} dengan total tabel,
\fashN{}, sehingga diperoleh proporsi pada
Gambar~\ref{photoClassifyProbTable}.
Distribusi peluang bersama variabel \var{mach\us{}learn}
dan \var{truth} ditampilkan pada
Gambar~\ref{photoClassifyDistribution}.

\begin{figure}[h]
\centering
\begin{tabular}{l rr r}
\hline
& \var{truth}: \resp{fashion} &
    \var{truth}: \resp{not} & Total  \\
\hline
\var{mach\us{}learn}: \resp{pred\us{}fashion} \hspace{0.5cm}
    & 0.1081 & 0.0121 & 0.1202 \\
\var{mach\us{}learn}: \resp{pred\us{}not}
    & 0.0615 & 0.8183 & 0.8798  \\
\hline
Total & 0.1696 & 0.8304 & 1.00 \\
\hline
\end{tabular}
\caption{Tabel peluang yang merangkum himpunan data
    \var{photo\us{}classify}.}
\label{photoClassifyProbTable}
\end{figure}

\begin{figure}[h]
\centering
\begin{tabular}{l c}
  \hline
Hasil bersama & Peluang \\
  \hline
\var{mach\us{}learn} adalah \resp{pred\us{}fashion}
    dan \var{truth} adalah \resp{fashion} & 0.1081 \\
\var{mach\us{}learn} adalah \resp{pred\us{}fashion}
    dan \var{truth} adalah \resp{not} & 0.0121 \\
\var{mach\us{}learn} adalah \resp{pred\us{}not}
    dan \var{truth} adalah \resp{fashion} & 0.0615 \\
\var{mach\us{}learn} adalah \resp{pred\us{}not}
    dan \var{truth} adalah \resp{not} & 0.8183 \\
   \hline
Total & 1.0000 \\
\hline
\end{tabular}
\caption{Distribusi peluang bersama untuk himpunan data \data{photo\us{}classify}.}
\label{photoClassifyDistribution}
\end{figure}

\D{\newpage}

\begin{exercisewrap}
\begin{nexercise}
Verifikasikan bahwa Gambar~\ref{photoClassifyDistribution} menyatakan
suatu distribusi probabilitas: peristiwa-peristiwanya saling lepas,
semua peluangnya tak negatif, dan jumlah seluruh peluangnya
adalah~1.\footnotemark
\end{nexercise}
\end{exercisewrap}
\footnotetext{Keempat kombinasi hasil tersebut saling lepas,
   semua peluangnya memang tak negatif, dan jumlah
   peluangnya adalah $0.1081 + 0.0121 + 0.0615 + 0.8183 = 1.00$.}

Peluang marginal dapat dihitung dari peluang bersama.
Misalnya, peluang foto bertema mode adalah jumlah peluang bagi
hasil-hasil dengan \var{truth} bernilai \resp{fashion}:%
\index{marginal probability|)}\index{joint probability|)}
\newcommand{\ultruthfashion}[0]
    {\underline{\var{truth} adalah \resp{fashion}}}%
\begin{align*}
P(\text{\ultruthfashion{}})
  &= P(\text{\var{mach\us{}learn} adalah \resp{pred\us{}fashion}
        dan \ultruthfashion{}}) \\
    & \qquad + P(\text{\var{mach\us{}learn} adalah \resp{pred\us{}not}
        dan \ultruthfashion{}}) \\
  &= 0.1081 + 0.0615 \\
  &= 0.1696
\end{align*}


\subsection{Mendefinisikan peluang bersyarat}

\index{conditional probability|(}

Pengklasifikasi ML memprediksi apakah sebuah foto berkaitan dengan mode,
meskipun hasilnya tidak sempurna.
Kita ingin memahami dengan lebih baik cara menggunakan informasi
dari variabel seperti \var{mach\us{}learn} untuk memperbaiki
perkiraan peluang bagi variabel kedua, yang dalam contoh ini
adalah \var{truth}.

Peluang bahwa sebuah foto acak dari himpunan data berkaitan dengan
mode adalah sekitar 0,17.
Jika kita mengetahui bahwa pengklasifikasi pembelajaran mesin memprediksi
foto tersebut berkaitan dengan mode, dapatkah kita memperoleh perkiraan yang lebih baik
atas peluang bahwa foto itu memang berkaitan dengan mode?
Tentu saja.
Untuk melakukannya, kita membatasi perhatian hanya pada \fashYA{} kasus
ketika pengklasifikasi ML memprediksi bahwa foto tersebut berkaitan dengan
mode, lalu melihat fraksi foto yang memang berkaitan
dengan mode:
\begin{align*}
P(\text{\var{truth} adalah \resp{fashion} jika
    \var{mach\us{}learn} adalah \resp{pred\us{}fashion}})
  = \frac{\fashYY{}}{\fashYA{}}
  = 0.900
\end{align*}
Kita menyebutnya \term{peluang bersyarat} karena
kita menghitung peluang di bawah suatu syarat:
pengklasifikasi ML memprediksi bahwa foto tersebut berkaitan dengan mode.

Peluang bersyarat terdiri atas dua bagian,
yaitu \term{hasil yang menjadi perhatian} dan \term{syarat}.
Syarat dapat dipandang sebagai informasi yang kita ketahui
benar, dan informasi ini biasanya dapat dinyatakan sebagai
hasil atau peristiwa yang diketahui.
Dalam notasi peluang, kita umumnya memisahkan
hasil yang menjadi perhatian dan syaratnya dengan
garis vertikal:
\begin{align*}
&& P(\text{\var{truth} adalah \resp{fashion} jika
    \var{mach\us{}learn} adalah \resp{pred\us{}fashion}}) \\
&& \quad = P(\text{\var{truth} adalah \resp{fashion}\ }|
    \text{\ \var{mach\us{}learn} adalah \resp{pred\us{}fashion}})
  = \frac{\fashYY{}}{\fashYA{}}
  = 0.900
\end{align*}
Garis vertikal ``$|$'' dibaca \emph{dengan syarat}.

\D{\newpage}

Pada persamaan terakhir, kita menghitung peluang sebuah foto
berkaitan dengan mode dengan syarat algoritma ML
memprediksi foto itu berkaitan dengan mode, sebagai fraksi:
\begin{align*}
& P(\text{\var{truth} adalah \resp{fashion}\ }|
    \text{\ \var{mach\us{}learn} adalah \resp{pred\us{}fashion}}) \\
  &\quad = \frac{\text{\# kasus ketika \var{truth} adalah \resp{fashion}
       dan \var{mach\us{}learn} adalah \resp{pred\us{}fashion}}}
     {\text{\# kasus ketika \var{mach\us{}learn} adalah \resp{pred\us{}fashion}}} \\
  &\quad = \frac{\fashYY{}}{\fashYA{}}
      = 0.900
\end{align*}
Kita hanya mempertimbangkan kasus yang memenuhi syarat,
\var{mach\us{}learn} bernilai \resp{pred\us{}fashion}, lalu
menghitung rasio kasus yang memenuhi
hasil yang menjadi perhatian, yaitu foto tersebut memang berkaitan dengan mode.

Sering kali peluang marginal dan bersama diberikan
sebagai pengganti data frekuensi.
Sebagai contoh, tingkat penyakit lazim disajikan dalam persentase,
bukan dalam bentuk frekuensi.
Kita ingin dapat menghitung peluang bersyarat
meskipun tidak tersedia frekuensi, dan kita menggunakan persamaan terakhir
sebagai pola untuk memahami teknik ini.

Kita hanya mempertimbangkan kasus yang memenuhi syarat,
yaitu algoritma ML memprediksi mode.
Di antara kasus-kasus ini, peluang bersyarat adalah
fraksi yang menyatakan hasil yang menjadi perhatian, yakni
foto tersebut berkaitan dengan mode.
Misalkan kita hanya diberi informasi dalam
Gambar~\ref{photoClassifyProbTable}, yaitu hanya data peluang.
Jika kita kemudian mengambil sampel 1.000 foto, kita memperkirakan
sekitar 12,0\% atau $0.120\times 1000 = 120$ foto akan diprediksi
berkaitan dengan mode (\var{mach\us{}learn} adalah \resp{pred\us{}fashion}).
Demikian pula, kita memperkirakan sekitar 10,8\% atau
$0.108\times 1000 = 108$ foto memenuhi kriteria informasi tersebut
sekaligus mewakili hasil yang menjadi perhatian.
Peluang bersyarat lalu dapat dihitung sebagai
\begin{align*}
&P(\text{\var{truth} adalah \resp{fashion}}\ |\ 
    \text{\var{mach\us{}learn} adalah \resp{pred\us{}fashion}}) \\
  &= \frac{\text{\# (\var{truth} adalah \resp{fashion}
      dan \var{mach\us{}learn} adalah \resp{pred\us{}fashion})}}
    {\text{\# (\var{mach\us{}learn} adalah \resp{pred\us{}fashion})}} \\
  &= \frac{108}{120}
		= \frac{0.108}{0.120}
		= 0.90
\end{align*}
Di sini kita menelaah tepatnya fraksi dari dua peluang,
0,108 dan 0,120, yang dapat kita tulis sebagai
\begin{align*}
&P(\text{\var{truth} adalah \resp{fashion} dan
    \var{mach\us{}learn} adalah \resp{pred\us{}fashion}})
\\[-1mm]
&\qquad\text{dan}\qquad
P(\text{\var{mach\us{}learn} adalah \resp{pred\us{}fashion}}).
\end{align*}
Fraksi kedua peluang ini merupakan contoh
rumus umum peluang bersyarat.

\begin{onebox}{Peluang bersyarat}
  Peluang bersyarat hasil $A$
  dengan syarat $B$ dihitung sebagai berikut:
  \begin{align*}
  P(A | B) = \frac{P(A\text{ dan }B)}{P(B)}
  \end{align*}
\end{onebox}

%\D{\newpage}

\begin{exercisewrap}
\begin{nexercise}
\label{fashionProbOfMLNotGivenTruthNot}%
(a) Tuliskan pernyataan berikut dalam notasi
peluang bersyarat:
``\emph{Peluang bahwa prediksi ML benar,
jika foto tersebut berkaitan dengan mode}''.
Di sini, syaratnya didasarkan pada status
\var{truth} foto, bukan pada algoritma ML. \\[1mm]
(b)~Tentukan peluang pada bagian (a).
Tabel~\vref{photoClassifyProbTable} mungkin membantu.\footnotemark
\end{nexercise}
\end{exercisewrap}
\footnotetext{(a) Jika foto tersebut berkaitan dengan mode dan
  prediksi algoritma ML benar, maka variabel keluaran algoritma ML
  harus bernilai \resp{pred\us{}fashion}:
  \begin{align*}
  P(\text{\var{mach\us{}learn} adalah \resp{pred\us{}fashion}}\ |
      \ \text{\var{truth} adalah \resp{fashion}})
  \end{align*}
  (b)~Persamaan peluang bersyarat menunjukkan bahwa kita
  harus terlebih dahulu mencari \\
  $P(\text{\var{mach\us{}learn} adalah \resp{pred\us{}fashion}
    dan \var{truth} adalah \resp{fashion}}) = 0.1081$
  dan $P(\text{\var{truth} adalah \resp{fashion}}) = 0.1696$. \\
  Rasionya kemudian menyatakan peluang bersyarat:
  $0.1081 / 0.1696 = 0.6374$.}

\begin{exercisewrap}
\begin{nexercise}
\label{whyCondProbSumTo1}%
(a)~Tentukan peluang bahwa algoritma tersebut keliru
jika diketahui foto itu berkaitan dengan mode. \\[1mm]
(b)~Dengan menggunakan jawaban bagian~(a) dan
Latihan Terbimbing~\ref{fashionProbOfMLNotGivenTruthNot}(b),
hitung
\begin{align*}
&P(\text{\var{mach\us{}learn} adalah \resp{pred\us{}fashion}}
    \ |\ \text{\var{truth} adalah \resp{fashion}}) \\
&\qquad  +\ P(\text{\var{mach\us{}learn} adalah \resp{pred\us{}not}}
    \ |\ \text{\var{truth} adalah \resp{fashion}})
\end{align*}
(c)~Berikan alasan intuitif yang menjelaskan mengapa jumlah
pada~(b) bernilai~1.\footnotemark
\end{nexercise}
\end{exercisewrap}
\footnotetext{(a)~Peluang ini adalah
  $\frac{P(\text{\var{mach\us{}learn} adalah \resp{pred\us{}not},
      \var{truth} adalah \resp{fashion}})}
    {P(\text{\var{truth} adalah \resp{fashion}})}
  = \frac{0.0615}{0.1696} = 0.3626$.
  (b)~Jumlahnya sama dengan~1.
  (c)~Dengan syarat foto tersebut berkaitan dengan mode,
      algoritma ML pasti memprediksinya berkaitan dengan mode
      atau tidak berkaitan dengan mode.
      Aturan komplemen tetap berlaku untuk peluang bersyarat,
      selama peluang-peluang tersebut dikondisikan pada
      informasi yang sama.}

\index{conditional probability|)}
\index{data!photo\_classify|)}


\subsection{Cacar di Boston, 1721}

\index{data!smallpox|(}

Himpunan data \data{smallpox} memuat sampel 6.224 orang dari tahun 1721 yang terpapar cacar di Boston.
Para dokter pada masa itu meyakini bahwa inokulasi, yakni memaparkan seseorang pada penyakit secara terkendali, dapat mengurangi peluang kematian.

Setiap kasus mewakili satu orang dengan dua variabel: \var{inoculated} dan \var{result}. Variabel \var{inoculated} memiliki dua tingkat, yaitu \resp{yes} atau \resp{no}, yang menunjukkan apakah orang tersebut diinokulasi. Variabel \var{result} memiliki hasil \resp{lived} atau \resp{died}. Data ini dirangkum dalam Tabel~\ref{smallpoxContingencyTable} dan~\ref{smallpoxProbabilityTable}.

\begin{figure}[h]
\centering
\begin{tabular}{ll rr r}
& & \multicolumn{2}{c}{diinokulasi} & \\
\cline{3-4}
& & \resp{yes} & \resp{no} & Total  \\
\cline{2-5}
		& \resp{lived}     & 238 & 5136 & 5374 \\
\raisebox{1.5ex}[0pt]{\var{result}} &  \resp{died} \hspace{0.5cm} & 6 & 844 & 850  \\
\cline{2-5}
	& Total & 244 & 5980 & 6224 \\
\end{tabular}
\caption{Tabel kontingensi untuk himpunan data \data{smallpox}.}
\label{smallpoxContingencyTable}
\end{figure}

\begin{figure}[h]
\centering
\begin{tabular}{ll rr r}
& & \multicolumn{2}{c}{diinokulasi} & \\
\cline{3-4}
& & \resp{yes} & \resp{no} & Total  \\
   \cline{2-5}
 & \resp{lived}     & 0.0382 & 0.8252 & 0.8634 \\
\raisebox{1.5ex}[0pt]{\var{result}} & \resp{died} \hspace{0.5cm} & 0.0010 & 0.1356  & 0.1366  \\
   \cline{2-5}
& Total & 0.0392 & 0.9608 & 1.0000 \\
\end{tabular}
\caption{Proporsi tabel untuk data \data{smallpox}, yang dihitung dengan membagi setiap frekuensi dengan total tabel, 6.224.}
\label{smallpoxProbabilityTable}
\end{figure}

%\D{\newpage}

\begin{exercisewrap}
\begin{nexercise} \label{probDiedIfNotInoculated}
Tuliskan, dalam notasi formal, peluang bahwa seseorang yang dipilih secara acak dan tidak diinokulasi meninggal akibat cacar, lalu hitung \mbox{peluang tersebut.}\footnotemark
\end{nexercise}
\end{exercisewrap}
\footnotetext{$P($\var{result} = \resp{died} $|$ \var{inoculated} = \resp{no}$) = \frac{P(\text{\var{result} = \resp{died} dan \var{inoculated} = \resp{no}})}{P(\text{\var{inoculated} = \resp{no}})} = \frac{0.1356}{0.9608} = 0.1411$.}

\begin{exercisewrap}
\begin{nexercise}
Tentukan peluang bahwa seseorang yang diinokulasi meninggal akibat cacar. Bagaimana hasil ini dibandingkan dengan hasil Latihan Terbimbing~\ref{probDiedIfNotInoculated}?\footnotemark
\end{nexercise}
\end{exercisewrap}
\footnotetext{$P($\var{result} = \resp{died} $|$ \var{inoculated} = \resp{yes}$) = \frac{P(\text{\var{result} = \resp{died} dan \var{inoculated} = \resp{yes}})}{P(\text{\var{inoculated} = \resp{yes}})} = \frac{0.0010}{0.0392} = 0.0255$ (jika kesalahan pembulatan dihindari, diperoleh $6 / 244 = 0.0246$). Tingkat kematian orang yang diinokulasi hanya sekitar 1~dari~40, sedangkan pada orang yang tidak diinokulasi sekitar 1~dari~7.}

\begin{exercisewrap}
\begin{nexercise}\label{SmallpoxInoculationObsExpExercise}
Penduduk Boston memilih sendiri apakah mereka akan diinokulasi. (a) Apakah penelitian ini bersifat observasional atau merupakan eksperimen? (b) Dapatkah kita menyimpulkan hubungan sebab-akibat dari data ini? (c) Apa saja variabel perancu potensial yang mungkin memengaruhi apakah seseorang \resp{lived} atau \resp{died}, sekaligus memengaruhi apakah orang tersebut diinokulasi?\footnotemark
\end{nexercise}
\end{exercisewrap}
\footnotetext{Jawaban singkat: (a)~Observasional. (b)~Tidak, kita tidak dapat menyimpulkan sebab-akibat dari penelitian observasional ini. (c)~Akses terhadap perawatan medis terbaru dan terbaik. Ada jawaban lain yang juga sah untuk bagian~(c).}


\subsection{Aturan perkalian umum}

Bagian~\ref{probabilityIndependence} memperkenalkan Aturan Perkalian untuk proses yang saling independen. Di sini kita menyajikan \term{Aturan Perkalian Umum} untuk peristiwa yang mungkin tidak saling independen.

\begin{onebox}{Aturan Perkalian Umum}
Jika $A$ dan $B$ menyatakan dua hasil atau peristiwa, maka \vspace{-1.5mm}
\begin{align*}
P(A\text{ dan }B) = P(A | B)\times P(B)
\end{align*} \vspace{-6.5mm} \par
$A$ dapat dipandang sebagai hasil yang menjadi perhatian dan $B$ sebagai syaratnya.
\end{onebox}

\noindent%
Aturan Perkalian Umum ini hanyalah bentuk susunan ulang dari persamaan peluang bersyarat.

%\D{\newpage}

\begin{examplewrap}
\begin{nexample}{Perhatikan himpunan data \data{smallpox}. Misalkan kita hanya diberi dua informasi: 96,08\% penduduk tidak diinokulasi, dan 85,88\% penduduk yang tidak diinokulasi bertahan hidup. Bagaimana kita dapat menghitung peluang bahwa seorang penduduk tidak diinokulasi dan tetap hidup?}
Kita akan menghitung jawabannya menggunakan Aturan Perkalian Umum, lalu memverifikasinya dengan Gambar~\ref{smallpoxProbabilityTable}. Kita ingin menentukan
\begin{align*}
P(\text{\var{result}
    = \resp{lived} dan \var{inoculated} = \resp{no}})
\end{align*}
dan diketahui bahwa
\begin{align*}
P(\text{\var{result}
    = \resp{lived} }|\text{ \var{inoculated} = \resp{no}})
    &= 0.8588 %\\
&&P(\text{\var{inoculated} = \resp{no}})
    = 0.9608
\end{align*}
Di antara 96,08\% orang yang tidak diinokulasi, 85,88\% bertahan hidup:
\begin{align*}
P(\text{\var{result} = \resp{lived}
        dan \var{inoculated} = \resp{no}})
    = 0.8588 \times 0.9608
    = 0.8251
\end{align*}
Perhitungan ini setara dengan Aturan Perkalian Umum. Kita dapat mengonfirmasi peluang ini pada Gambar~\ref{smallpoxProbabilityTable}, di perpotongan \resp{no} dan \resp{lived} (dengan sedikit kesalahan pembulatan).
\end{nexample}
\end{examplewrap}

\begin{exercisewrap}
\begin{nexercise}
Gunakan $P($\var{inoculated} = \resp{yes}$) = 0.0392$ dan $P($\var{result} = \resp{lived} $|$ \var{inoculated} = \resp{yes}$) = 0.9754$ untuk menentukan peluang bahwa seseorang diinokulasi dan tetap hidup.\footnotemark
\end{nexercise}
\end{exercisewrap}
\footnotetext{Jawabannya adalah 0,0382, yang dapat diverifikasi menggunakan Gambar~\ref{smallpoxProbabilityTable}.}

%\D{\newpage}

\begin{exercisewrap}
\begin{nexercise}
Jika 97,54\% orang yang diinokulasi tetap hidup,
berapa proporsi orang yang diinokulasi yang meninggal?\footnotemark{}
\end{nexercise}
\end{exercisewrap}
\footnotetext{Hanya ada dua hasil yang mungkin: \resp{lived} atau \resp{died}. Artinya, 100\% - 97,54\% = 2,46\% orang yang diinokulasi meninggal.}

\begin{onebox}{Jumlah peluang bersyarat}
Misalkan $A_1$, ..., $A_k$ menyatakan semua hasil yang saling lepas bagi suatu variabel atau proses. Jika $B$ adalah sebuah peristiwa, yang mungkin berasal dari variabel atau proses lain, maka: \vspace{-1mm}
\begin{align*}
P(A_1|B) + \cdots + P(A_k|B) = 1
\end{align*}%
\vspace{-5.5mm} \par
Aturan komplemen juga berlaku apabila sebuah peristiwa dan komplemennya dikondisikan pada informasi yang sama: \vspace{-1.5mm}
\begin{align*}
P(A | B) = 1 - P(A^c | B)
\end{align*}
\end{onebox}

\begin{exercisewrap}
\begin{nexercise}
Berdasarkan peluang yang dihitung di atas, apakah inokulasi tampaknya efektif mengurangi risiko kematian akibat cacar?\footnotemark
\end{nexercise}
\end{exercisewrap}
\footnotetext{Ukuran sampelnya besar jika dibandingkan dengan perbedaan tingkat kematian antara kelompok ``diinokulasi'' dan ``tidak diinokulasi'', sehingga tampaknya terdapat hubungan antara \var{inoculated} dan \var{result}. Namun, sebagaimana dicatat dalam jawaban Latihan Terbimbing~\ref{SmallpoxInoculationObsExpExercise}, penelitian ini bersifat observasional sehingga kita tidak dapat memastikan adanya hubungan sebab-akibat. (Penelitian lebih lanjut menunjukkan bahwa inokulasi efektif menurunkan tingkat kematian.)}


%\D{\newpage}

\subsection{Pertimbangan independensi dalam peluang bersyarat}

Jika dua peristiwa saling independen, mengetahui hasil salah satunya tidak memberikan informasi tentang peristiwa lainnya. Dengan peluang bersyarat, kita dapat menunjukkan bahwa pernyataan ini benar secara matematis.

\begin{exercisewrap}
\begin{nexercise} \label{condProbOfRollingA1AfterOne1}
Misalkan $X$ dan $Y$ menyatakan hasil pelemparan dua dadu.\footnotemark
\begin{enumerate}[(a)]
\setlength{\itemsep}{0mm}
\item Berapa peluang bahwa dadu pertama, $X$, menghasilkan \resp{1}?
\item Berapa peluang bahwa $X$ dan $Y$ sama-sama menghasilkan \resp{1}?
\item Gunakan rumus peluang bersyarat untuk menghitung $P(Y =$ \resp{1}$\ |\ X = $ \resp{1}$)$.
\item Berapa nilai $P(Y=1)$? Apakah hasil ini berbeda dari jawaban bagian (c)? Jelaskan.
\end{enumerate}
\end{nexercise}
\end{exercisewrap}
\footnotetext{Jawaban singkat: (a) $1/6$. (b) $1/36$. (c)~$\frac{P(Y = \text{ \resp{1} dan }X=\text{ \resp{1}})}{P(X=\text{ \resp{1}})} = \frac{1/36}{1/6} = 1/6$. (d)~Peluangnya sama seperti pada bagian~(c): $P(Y=1)=1/6$. Peluang bahwa $Y=1$ tidak berubah setelah kita mengetahui $X$, sesuai dengan fakta bahwa $X$ dan $Y$ saling independen.}

Pada Latihan Terbimbing~\ref{condProbOfRollingA1AfterOne1}(c), kita dapat menunjukkan bahwa informasi yang menjadi syarat tidak berpengaruh dengan menggunakan Aturan Perkalian untuk proses yang saling independen:
\begin{align*}
P(Y=\text{\resp{1}}\ |\ X=\text{\resp{1}})
    &= \frac{P(Y=\text{\resp{1} dan }X=\text{\resp{1}})}
      {P(X=\text{\resp{1}})} \\
    &= \frac{P(Y=\text{\resp{1}}) \times
        \color{oiGB}P(X=\text{\resp{1}})}
      {\color{oiGB}P(X=\text{\resp{1}})} \\
    &= P(Y=\text{\resp{1}}) \\
\end{align*}

\begin{exercisewrap}
\begin{nexercise}
Ron sedang mengamati meja rolet di kasino dan melihat bahwa lima hasil terakhir adalah \resp{black}. Ia beranggapan bahwa peluang memperoleh \resp{black} enam kali berturut-turut sangat kecil (sekitar $1/64$), lalu mempertaruhkan seluruh gajinya pada merah. Apa yang keliru dalam penalarannya?\footnotemark
\end{nexercise}
\end{exercisewrap}
\footnotetext{Ia lupa bahwa putaran rolet berikutnya independen dari putaran-putaran sebelumnya. Kasino memang menggunakan praktik ini: mereka menampilkan beberapa hasil terakhir dari berbagai permainan taruhan untuk mengecoh penjudi yang tidak waspada agar percaya bahwa peluang berpihak kepada mereka. Kekeliruan ini disebut \term{sesat pikir penjudi}.}


\D{\newpage}

\subsection{Diagram pohon}

\index{data!smallpox|)}
\index{tree diagram|(}

\termsub{Diagram pohon}{diagram pohon} merupakan alat untuk menata hasil dan peluang berdasarkan struktur data. Diagram ini paling berguna ketika dua proses atau lebih terjadi secara berurutan dan setiap proses dikondisikan pada proses-proses sebelumnya.

Data \data{smallpox} sesuai dengan gambaran ini. Populasi terbagi menurut \var{inoculation}: \resp{yes} dan \resp{no}. Setelah pembagian ini, tingkat kelangsungan hidup diamati pada setiap kelompok. Struktur tersebut tercermin dalam \term{diagram pohon} pada Gambar~\ref{smallpoxTreeDiagram}. Cabang pertama untuk \var{inoculation} disebut \term{cabang utama}, sedangkan cabang-cabang lainnya disebut \term{cabang sekunder}.

\begin{figure}[ht]
\centering
\Figure[Diagram pohon dengan cabang utama "Diinokulasi" dan cabang sekunder "Hasil". Cabang utama Diinokulasi menuju dua pilihan: "Ya" dengan peluang 0,0392 dan "Tidak" dengan peluang 0,9608. Masing-masing cabang ini memiliki cabang sekunder berisi peluang "Hasil" yang dikondisikan pada "Diinokulasi". Cabang Diinokulasi--Ya terbagi menjadi "Hidup" (0,9754) dan "Meninggal" (0,0246). Cabang-cabang ini juga menampilkan hasil perkalian peluang sepanjang cabang. Sebagai contoh, cabang Ya-dan-Hidup mengalikan 0,0392 dengan 0,9754 dan menghasilkan 0,03824. Cabang Ya-dan-Meninggal memiliki peluang hasil perkalian 0,00096. Selanjutnya, cabang utama "Tidak" juga memiliki cabang sekunder Hidup dan Meninggal, masing-masing dengan peluang bersyarat 0,8589 dan 0,1411. Hasil perkalian peluang pada setiap rangkaian cabang juga ditampilkan: Tidak-dan-Hidup sebesar 0,82523 dan Tidak-dan-Meninggal sebesar 0,13557.]{0.93}{smallpoxTreeDiagram}
\caption{Diagram pohon untuk himpunan data \data{smallpox}.}
\label{smallpoxTreeDiagram}
\end{figure}

Diagram pohon diberi keterangan peluang marginal dan peluang bersyarat, seperti pada Gambar~\ref{smallpoxTreeDiagram}. Diagram ini membagi data cacar menurut \var{inoculation} menjadi kelompok \resp{yes} dan \resp{no}, dengan peluang marginal masing-masing 0,0392 dan 0,9608. Cabang-cabang sekunder dikondisikan pada cabang pertama, sehingga kita memberikan peluang bersyarat pada cabang-cabang ini. Sebagai contoh, cabang teratas pada Gambar~\ref{smallpoxTreeDiagram} adalah peluang bahwa \var{result} = \resp{lived} dengan syarat \var{inoculated} = \resp{yes}. Kita dapat (dan biasanya memang) menghitung peluang bersama di ujung setiap cabang dengan mengalikan bilangan-bilangan yang ditemui saat bergerak dari kiri ke kanan. Peluang bersama ini dihitung menggunakan Aturan Perkalian Umum:
\begin{align*}
& P(\text{\var{inoculated} = \resp{yes}
    dan \var{result} = \resp{lived}}) \\
  &\quad = P(\text{\var{inoculated} = \resp{yes}})\times
      P(\text{\var{result} = \resp{lived}}|
          \text{\var{inoculated} = \resp{yes}}) \\
  &\quad = 0.0392\times 0.9754=0.0382
\end{align*}

\begin{examplewrap}
\begin{nexample}{Perhatikan ujian tengah semester dan ujian akhir pada sebuah kelas statistika. Misalkan 13\% mahasiswa memperoleh \resp{A} pada ujian tengah semester. Di antara mahasiswa yang memperoleh \resp{A} pada ujian tengah semester, 47\% memperoleh \resp{A} pada ujian akhir; sementara itu, 11\% mahasiswa yang memperoleh nilai di bawah \resp{A} pada ujian tengah semester memperoleh \resp{A} pada ujian akhir. Anda mengambil secara acak satu lembar ujian akhir dan melihat bahwa mahasiswa tersebut memperoleh \resp{A}. Berapa peluang mahasiswa itu juga memperoleh \resp{A} pada ujian tengah semester?} \label{exerciseForTreeDiagramOfStudentGettingAOnMidtermGivenThatSheGotAOnFinal}
Tujuan akhirnya adalah mencari $P(\text{\var{midterm} = \resp{A}} | \text{\var{final} = \resp{A}})$. Untuk menghitung peluang bersyarat ini, kita memerlukan peluang berikut:
\begin{align*}
P(\text{\var{midterm} = \resp{A} dan \var{final} = \resp{A}})
  \qquad\text{dan}\qquad
  P(\text{\var{final} = \resp{A}})
\end{align*}
Namun, informasi ini tidak diberikan, dan cara menghitung peluang tersebut tidak langsung terlihat. Karena langkah berikutnya belum jelas, informasi ini sebaiknya kita tata dalam sebuah diagram pohon:
\begin{center}
\Figure[Diagram pohon dengan cabang utama "Ujian Tengah Semester" dan cabang sekunder "Ujian Akhir". Cabang utama Ujian Tengah Semester menuju dua pilihan: "A" dengan peluang 0,13 dan "Lainnya" dengan peluang 0,87. Masing-masing cabang memiliki cabang sekunder berisi peluang bersyarat "Ujian Akhir" dengan syarat "Ujian Tengah Semester". Cabang Ujian-Tengah-Semester--A terbagi menjadi cabang "A", dengan peluang bersyarat 0,47 dan peluang akhir A-dan-A sebesar 0,0611, serta cabang sekunder "lainnya", dengan peluang bersyarat 0,53 dan peluang akhir A-dan-lainnya sebesar 0,0689. Selanjutnya, cabang utama Ujian-Tengah-Semester--Lainnya juga memiliki cabang sekunder Ujian-Akhir--A dengan peluang bersyarat 0,11 dan peluang akhir 0,0957, serta cabang Ujian-Akhir--lainnya dengan peluang bersyarat 0,89 dan peluang akhir 0,7743.]{0.85}{testTree}
\end{center}
Ketika menyusun diagram pohon, variabel yang dilengkapi peluang marginal sering digunakan untuk membentuk cabang utama; dalam kasus ini, peluang marginal diberikan untuk nilai ujian tengah semester. Nilai ujian akhir, yang bersesuaian dengan peluang bersyarat yang diberikan, ditampilkan pada cabang sekunder.

Setelah diagram pohon tersusun, kita dapat menghitung peluang yang diperlukan:
\begin{align*}
&P(\text{\var{midterm} = \resp{A} dan \var{final} = \resp{A}}) = 0.0611 \\
&P(\text{\underline{\color{black}\var{final} = \resp{A}}})  \\
& \quad= P(\text{\var{midterm} = \resp{other} dan \underline{\color{black}\var{final} = \resp{A}}}) + P(\text{\var{midterm} = \resp{A} dan \underline{\color{black}\var{final} = \resp{A}}}) \\
& \quad= 0.0957 + 0.0611  = 0.1568
\end{align*}
Peluang marginal $P($\var{final} = \resp{A}$)$ dihitung dengan menjumlahkan semua peluang bersama di sisi kanan pohon yang bersesuaian dengan \var{final} = \resp{A}. Sekarang kita dapat mengambil rasio kedua peluang tersebut:
\begin{align*}
P(\text{\var{midterm} = \resp{A}} | \text{\var{final} = \resp{A}})
  &= \frac{P(\text{\var{midterm} = \resp{A}
      dan \var{final} = \resp{A}})}
    {P(\text{\var{final} = \resp{A}})} \\
  &= \frac{0.0611}{0.1568} = 0.3897
\end{align*}
Peluang bahwa mahasiswa tersebut juga memperoleh A pada ujian tengah semester adalah sekitar 0,39.
\end{nexample}
\end{examplewrap}

%\begin{figure}[ht]
%  \centering
%  \Figure{0.9}{testTree}
%  \caption{Diagram pohon yang menggambarkan variabel \var{midterm}
%      dan \var{final}.}
%  \label{testTree}
%\end{figure}

\begin{exercisewrap}
\begin{nexercise}
Setelah mengikuti mata kuliah statistika pengantar, 78\% mahasiswa dapat menyusun diagram pohon dengan benar. Di antara mahasiswa yang dapat menyusun diagram pohon, 97\% lulus, sedangkan hanya 57\% mahasiswa yang tidak dapat menyusun diagram pohon yang lulus. (a)~Susun informasi ini dalam sebuah diagram pohon. (b)~Berapa peluang bahwa mahasiswa yang dipilih secara acak lulus? (c)~Hitung peluang bahwa seorang mahasiswa mampu menyusun diagram pohon jika diketahui mahasiswa tersebut lulus.\footnotemark
\end{nexercise}
\end{exercisewrap}
\footnotetext{%\begin{minipage}[t]{0.27\linewidth}
(a) Diagram pohon ditampilkan di sebelah kanan. \\
(b)~Tentukan dua peluang bersama yang mewakili
    mahasiswa yang lulus, lalu jumlahkan:
    $P($lulus$) = 0.7566+0.1254= 0.8820$. \\
(c)~$P($mampu menyusun diagram pohon $|$ lulus$)
   = \frac{0.7566}{0.8820} = 0.8578$. \\ %\vspace{15mm} \\
%\end{minipage}
%\begin{minipage}[c]{0.7\linewidth}
\Figure[Diagram pohon dengan cabang utama "Mampu menyusun diagram pohon" dan cabang sekunder "Lulus mata kuliah". Cabang utama Mampu-menyusun-diagram-pohon menuju dua pilihan: "Ya" dengan peluang 0,78 dan "Tidak" dengan peluang 0,22. Masing-masing cabang memiliki cabang sekunder berisi peluang "Lulus Mata Kuliah" dengan syarat "Mampu menyusun diagram pohon". Cabang utama Ya terbagi menjadi cabang "Lulus", dengan peluang bersyarat 0,97 dan peluang akhir Ya-dan-Lulus sebesar 0,7566, serta cabang sekunder "Tidak lulus", dengan peluang bersyarat 0,03 dan peluang akhir Ya-dan-Tidak-Lulus sebesar 0,0234. Selanjutnya, cabang utama Tidak juga memiliki cabang sekunder Lulus dengan peluang bersyarat 0,57 dan peluang akhir 0,1254, serta cabang Tidak Lulus dengan peluang bersyarat 0,43 dan peluang akhir 0,0946.]{0.7}{treeDiagramAndPass}}% \vspace{-25mm}
%\end{minipage}}


\subsection{Teorema Bayes}
\label{bayesTheoremSubsection}

\index{Bayes' Theorem|(}

Dalam banyak kasus, kita diberi peluang bersyarat berbentuk
\begin{align*}
P(\text{pernyataan tentang variabel 1 } | \text{ pernyataan tentang variabel 2})
\end{align*}
padahal yang ingin kita ketahui adalah peluang bersyarat dengan arah terbalik:
\begin{align*}
P(\text{pernyataan tentang variabel 2 } | \text{ pernyataan tentang variabel 1})
\end{align*}
Diagram pohon dapat digunakan untuk mencari peluang bersyarat kedua jika peluang yang pertama diketahui. Namun, kadang-kadang skenarionya tidak dapat digambarkan dengan diagram pohon. Dalam kasus seperti ini, kita dapat menerapkan rumus umum yang sangat berguna: Teorema Bayes.

Pertama, kita telaah dengan cermat sebuah contoh pembalikan peluang bersyarat yang masih dapat diselesaikan menggunakan diagram pohon.

\D{\newpage}

\begin{examplewrap}
\begin{nexample}{Di Kanada, sekitar 0,35\% perempuan berusia di atas 40 tahun
    mengidap kanker payudara pada suatu tahun.
    Mamogram lazim digunakan untuk penapisan, tetapi tidak sempurna.
    Pada 11\% pasien yang mengidap kanker, hasilnya \term{negatif palsu}:
    pemeriksaan menyatakan tidak ada kanker padahal kanker ada.
    Pada 7\% pasien tanpa kanker, hasilnya \term{positif palsu}:
    pemeriksaan menyatakan ada kanker padahal tidak.
    Jika seorang perempuan di atas 40 tahun dipilih secara acak dan hasil
    mamogramnya positif, berapa peluang ia benar-benar mengidap kanker payudara?}

\label{probBreastCancerGivenPositiveTestExample}

Informasi ini langsung memberi peluang hasil pemeriksaan positif jika pasien mengidap kanker ($1.00-0.11=0.89$). Namun, kita mencari peluang dengan arah terbalik: peluang mengidap kanker dengan syarat hasil pemeriksaan positif. (Dalam istilah penapisan yang kurang intuitif, hasil mamogram \emph{positif} menunjukkan kemungkinan adanya kanker.) Peluang terbalik ini dapat diuraikan menjadi dua bagian:
\begin{align*}
P(\text{mengidap KP } | \text{ mamogram$^+$}) = \frac{P(\text{mengidap KP dan mamogram$^+$})}{P(\text{mamogram$^+$})}
\end{align*}
dengan ``mengidap KP'' sebagai singkatan bahwa pasien mengidap
kanker payudara dan ``mamogram$^+$'' berarti hasil penapisan mamogram
positif.
Kita dapat menyusun diagram pohon untuk peluang-peluang ini:
\begin{center}
\Figure[Diagram pohon dengan cabang utama "Kondisi sebenarnya" dan cabang sekunder "Mamogram". Cabang utama Kondisi sebenarnya menuju dua pilihan: "Kanker" dengan peluang 0,0035 dan "Tanpa Kanker" dengan peluang 0,9965. Masing-masing cabang memiliki cabang sekunder berisi peluang bersyarat untuk hasil mamogram "Positif" dan "Negatif", dengan syarat apakah kondisi sebenarnya berupa kanker atau bukan. Cabang utama Kanker terbagi menjadi cabang "Positif", dengan peluang bersyarat 0,89 dan peluang akhir Kanker-dan-Positif sebesar 0,00312, serta cabang sekunder "Negatif", dengan peluang bersyarat 0,11 dan peluang akhir Kanker-dan-Negatif sebesar 0,00038. Selanjutnya, cabang utama Tanpa-Kanker juga memiliki cabang sekunder Positif dengan peluang bersyarat 0,07 dan peluang akhir 0,06976, serta cabang Negatif dengan peluang bersyarat 0,93 dan peluang akhir 0,92675.]{0.9}{BreastCancerTreeDiagram}
\end{center}
Peluang bahwa pasien mengidap kanker payudara
dan hasil mamogramnya positif adalah
\begin{align*}
P(\text{mengidap KP dan mamogram$^+$}) &= P(\text{mamogram$^+$ } | \text{ mengidap KP})P(\text{mengidap KP}) \\
	&= 0.89\times 0.0035 = 0.00312
\end{align*}
Peluang hasil pemeriksaan positif adalah jumlah peluang dua skenario yang bersesuaian:
\begin{align*}
P(\text{\underline{\color{black}mamogram$^+$}})
  &= P(\text{\underline{\color{black}mamogram$^+$} dan mengidap KP}) \\
  &\qquad\qquad
      + P(\text{\underline{\color{black}mamogram$^+$} dan tidak mengidap KP})\\
  &= P(\text{mengidap KP})P(\text{mamogram$^+$ } | \text{ mengidap KP}) \\
  &\qquad\qquad
      + P(\text{tidak mengidap KP})P(\text{mamogram$^+$ } | \text{ tidak mengidap KP}) \\
  &= 0.0035\times 0.89 + 0.9965\times 0.07 = 0.07288
\end{align*}
Jadi, jika hasil penapisan mamogram seorang pasien positif, peluang bahwa pasien tersebut mengidap kanker payudara adalah
\begin{align*}
P(\text{mengidap KP } | \text{ mamogram$^+$})
	&= \frac{P(\text{mengidap KP dan mamogram$^+$})}{P(\text{mamogram$^+$})}\\
	&= \frac{0.00312}{0.07288} \approx 0.0428
\end{align*}
Dengan kata lain, meskipun hasil penapisan mamogram seorang pasien positif, peluang bahwa pasien tersebut mengidap kanker payudara tetap hanya sekitar~4\%.
\end{nexample}
\end{examplewrap}

%\begin{figure}[h]
%  \centering
%  \Figure{0.75}{BreastCancerTreeDiagram}
%  \caption{Diagram pohon untuk
%      Contoh~\ref{probBreastCancerGivenPositiveTestExample}.}%,
%      %menghitung peluang bahwa pasien acak dengan hasil mamogram
%      %positif benar-benar mengidap kanker payudara.}
%\label{BreastCancerTreeDiagram}
%\end{figure}

\D{\newpage}

Contoh~\ref{probBreastCancerGivenPositiveTestExample} menunjukkan mengapa dokter sering melakukan pemeriksaan lanjutan setelah hasil pemeriksaan pertama positif. Jika suatu kondisi medis jarang terjadi, satu hasil positif umumnya belum bersifat meyakinkan.

Perhatikan kembali persamaan terakhir pada Contoh~\ref{probBreastCancerGivenPositiveTestExample}.
Dari diagram pohon, kita melihat bahwa pembilang (bagian atas pecahan) sama dengan hasil kali berikut:
\begin{align*}
P(\text{mengidap KP dan mamogram$^+$}) = P(\text{mamogram$^+$ } | \text{ mengidap KP})P(\text{mengidap KP})
\end{align*}
Penyebutnya -- peluang bahwa hasil penapisan positif -- sama dengan jumlah peluang setiap skenario penapisan positif:
\begin{align*}
P(\text{\underline{\color{black}mamogram$^+$}})
	&= P(\text{\underline{\color{black}mamogram$^+$} dan tidak mengidap KP})
		+ P(\text{\underline{\color{black}mamogram$^+$} dan mengidap KP})
\end{align*}
Dalam contoh tersebut, diagram pohon digunakan untuk menguraikan setiap peluang di ruas kanan menjadi hasil kali peluang bersyarat dan peluang marginal.
\begin{align*}
P(\text{mamogram$^+$})
	&= P(\text{mamogram$^+$ dan tidak mengidap KP}) + P(\text{mamogram$^+$ dan mengidap KP}) \\
	&= P(\text{mamogram$^+$ } | \text{ tidak mengidap KP})P(\text{tidak mengidap KP}) \\
			   &\qquad\qquad + P(\text{mamogram$^+$ } | \text{ mengidap KP})P(\text{mengidap KP})
\end{align*}
Penerapan Teorema Bayes terlihat ketika bentuk-bentuk peluang yang dihasilkan disubstitusikan ke pembilang dan penyebut peluang bersyarat semula.
\begin{align*}
& P(\text{mengidap KP } | \text{ mamogram$^+$})  \\
& \qquad= \frac{P(\text{mamogram$^+$ } | \text{ mengidap KP})P(\text{mengidap KP})}
	{\begin{gathered}
	 P(\text{mamogram$^+$ } | \text{ tidak mengidap KP})P(\text{tidak mengidap KP}) \\
	 {}+P(\text{mamogram$^+$ } | \text{ mengidap KP})P(\text{mengidap KP})
	 \end{gathered}}
\end{align*}

\begin{onebox}{Teorema Bayes: membalik peluang}
Perhatikan peluang bersyarat berikut untuk variabel 1 dan variabel 2:\vspace{-1.5mm}
\begin{align*}
P(\text{hasil $A_1$ dari variabel 1 } | \text{ hasil $B$ dari variabel 2})
\end{align*}
Teorema Bayes menyatakan bahwa peluang bersyarat ini dapat ditentukan dengan pecahan berikut:\vspace{-1.5mm}
\begin{align*}
\frac{P(B | A_1) P(A_1)}
	{P(B | A_1) P(A_1) + P(B | A_2) P(A_2) + \cdots + P(B | A_k) P(A_k)}
\end{align*}
dengan $A_2$, $A_3$, ..., dan $A_k$ menyatakan semua hasil lain yang mungkin dari variabel pertama.\index{Bayes' Theorem|textbf}
\end{onebox}

Teorema Bayes merupakan generalisasi dari langkah yang telah kita lakukan
dengan diagram pohon.
Pembilangnya menunjukkan peluang memperoleh
$A_1$ dan~$B$ sekaligus.
Penyebutnya adalah peluang marginal memperoleh~$B$.
Bagian bawah pecahan ini tampak panjang dan
rumit karena kita harus menjumlahkan peluang dari semua cara berbeda untuk memperoleh $B$. Kita selalu melakukan langkah ini saat menggunakan diagram pohon. Namun, biasanya langkah tersebut dikerjakan secara terpisah sehingga tidak tampak serumit ini.

\noindent%
Untuk menerapkan Teorema Bayes dengan benar, ada dua langkah persiapan:
\begin{enumerate}
\setlength{\itemsep}{0mm}
\item[(1)] Pertama, tentukan peluang marginal setiap hasil yang mungkin dari variabel pertama: $P(A_1)$, $P(A_2)$, ..., $P(A_k)$.
\item[(2)] Kemudian, tentukan peluang hasil $B$ dengan syarat setiap skenario yang mungkin bagi variabel pertama: $P(B | A_1)$, $P(B | A_2)$, ..., $P(B | A_k)$.
\end{enumerate}
Setelah semua peluang ini ditentukan, nilainya dapat langsung dimasukkan ke dalam rumus.
Teorema Bayes biasanya merupakan pilihan yang baik ketika terdapat begitu banyak skenario sehingga diagram pohon akan rumit.

\begin{exercisewrap}
\begin{nexercise} \label{exerciseForParkingLotOnCampusBeingFullAndWhetherOrNotThereIsASportingEvent}
Jose mengunjungi kampus setiap Kamis malam. Namun, pada beberapa malam gedung parkir penuh, sering kali karena ada acara kampus. Acara akademik berlangsung pada 35\% malam, acara olahraga pada 20\% malam, dan tidak ada acara pada 45\% malam. Saat ada acara akademik, gedung parkir penuh sekitar 25\% dari waktu tersebut; saat ada acara olahraga, gedung parkir penuh pada 70\% malam. Pada malam tanpa acara, gedung parkir hanya penuh sekitar 5\% dari waktu tersebut. Jika Jose datang ke kampus dan mendapati gedung parkir penuh, berapa peluang bahwa sedang berlangsung acara olahraga? Gunakan diagram pohon untuk menyelesaikan masalah ini.\footnotemark
\end{nexercise}
\end{exercisewrap}
\footnotetext{\begin{minipage}[t]{0.27\linewidth}
Diagram pohon dengan tiga cabang utama ditampilkan di sebelah kanan. Selanjutnya, kita menentukan dua peluang dari diagram tersebut. (1) Peluang bahwa ada acara olahraga dan gedung parkir penuh: 0,14. (2) Peluang gedung parkir penuh: $0.0875 + 0.14 + 0.0225 = 0.25$. Jawabannya adalah rasio kedua peluang tersebut: $\frac{0.14}{0.25} = 0.56$. Jika gedung parkir penuh, peluang bahwa ada acara olahraga adalah 56\%. \vspace{0.1mm} \\\ 
\end{minipage}
\begin{minipage}[c]{0.65\linewidth}
\Figure[Diagram pohon dengan cabang utama "Acara" dan cabang sekunder "Gedung parkir penuh". Cabang utama "Acara" memiliki tiga kemungkinan: "Akademik" dengan peluang 0,35, "Olahraga" dengan peluang 0,20, dan "Tidak ada" dengan peluang 0,45. Masing-masing dari ketiga cabang ini memiliki dua cabang sekunder. Cabang utama "Akademik" terbagi menjadi cabang "Penuh" dengan peluang bersyarat 0,25 dan peluang akhir Akademik-dan-Penuh sebesar 0,0875, serta cabang sekunder "Tempat Tersedia" dengan peluang bersyarat 0,75 dan peluang akhir Akademik-dan-Tempat-Tersedia sebesar 0,2625. Cabang utama "Olahraga" terbagi menjadi cabang "Penuh" dengan peluang bersyarat 0,7 dan peluang akhir Olahraga-dan-Penuh sebesar 0,14, serta cabang sekunder "Tempat Tersedia" dengan peluang bersyarat 0,3 dan peluang akhir Olahraga-dan-Tempat-Tersedia sebesar 0,06. Cabang utama "Tidak ada" terbagi menjadi cabang "Penuh" dengan peluang bersyarat 0,05 dan peluang akhir Tidak-Ada-dan-Penuh sebesar 0,0225, serta cabang sekunder "Tempat Tersedia" dengan peluang bersyarat 0,95 dan peluang akhir Tidak-Ada-dan-Penuh sebesar 0,4275.]{}{treeDiagramGarage}\vspace{-45mm}
\end{minipage}}

\begin{examplewrap}
\begin{nexample}{Di sini kita menyelesaikan masalah yang sama seperti pada Latihan Terbimbing~\ref{exerciseForParkingLotOnCampusBeingFullAndWhetherOrNotThereIsASportingEvent}, tetapi kali ini menggunakan Teorema Bayes.}
Hasil yang menjadi perhatian adalah ada atau tidaknya acara olahraga (kita sebut $A_1$), sedangkan syaratnya adalah gedung parkir penuh ($B$). Misalkan $A_2$ menyatakan acara akademik dan $A_3$ menyatakan tidak adanya acara di kampus. Peluang yang diberikan dapat ditulis sebagai
\begin{align*}
&P(A_1) = 0.2 &&P(A_2) = 0.35 &&P(A_3) = 0.45 \\
&P(B | A_1) = 0.7 &&P(B | A_2) = 0.25 &&P(B | A_3) = 0.05
\end{align*}
Teorema Bayes dapat digunakan untuk menghitung peluang adanya acara olahraga ($A_1$) dengan syarat gedung parkir penuh ($B$):
\begin{align*}
P(A_1 | B) &= \frac{P(B | A_1) P(A_1)}{P(B | A_1) P(A_1) + P(B | A_2) P(A_2) + P(B | A_3) P(A_3)} \\
		&= \frac{(0.7)(0.2)}{(0.7)(0.2) + (0.25)(0.35) + (0.05)(0.45)} \\
		&= 0.56 
\end{align*}
Berdasarkan informasi bahwa gedung parkir penuh, peluang bahwa acara olahraga sedang berlangsung di kampus pada malam itu adalah 56\%.
\end{nexample}
\end{examplewrap}

\D{\newpage}

\begin{exercisewrap}
\begin{nexercise} \label{exerciseForParkingLotOnCampusBeingFullAndWhetherOrNotThereIsAnAcademicEvent}
Gunakan informasi pada latihan dan contoh sebelumnya untuk memverifikasi bahwa peluang adanya acara akademik dengan syarat gedung parkir penuh adalah 0,35.\footnotemark
\end{nexercise}
\end{exercisewrap}
\footnotetext{Jawaban singkat:
\begin{align*}
P(A_2 | B) &= \frac{P(B | A_2) P(A_2)}{P(B | A_1) P(A_1) + P(B | A_2) P(A_2) + P(B | A_3) P(A_3)} \\
		&= \frac{(0.25)(0.35)}{(0.7)(0.2) + (0.25)(0.35) + (0.05)(0.45)} \\
		&= 0.35
\end{align*}}

\begin{exercisewrap}
\begin{nexercise} \label{exerciseForParkingLotOnCampusBeingFullAndWhetherOrNotThereIsNoEvent}
Dalam Latihan Terbimbing~\ref{exerciseForParkingLotOnCampusBeingFullAndWhetherOrNotThereIsASportingEvent} dan~\ref{exerciseForParkingLotOnCampusBeingFullAndWhetherOrNotThereIsAnAcademicEvent}, Anda menemukan bahwa jika gedung parkir penuh, peluang adanya acara olahraga adalah 0,56 dan peluang adanya acara akademik adalah 0,35. Dengan informasi ini, hitung $P($tidak ada acara $|$ gedung parkir penuh$)$.\footnotemark
\end{nexercise}
\end{exercisewrap}
\footnotetext{Setiap peluang dikondisikan pada informasi yang sama, yaitu gedung parkir penuh, sehingga aturan komplemen dapat digunakan: $1.00 - 0.56 - 0.35 = 0.09$.}

Beberapa latihan terakhir memberikan cara untuk memperbarui keyakinan kita mengenai apakah sedang ada acara olahraga, acara akademik, atau tidak ada acara di kampus berdasarkan informasi bahwa gedung parkir penuh. Strategi \emph{memperbarui keyakinan} dengan Teorema Bayes ini merupakan landasan sebuah cabang statistika yang disebut \term{statistika Bayes}. Walaupun statistika Bayes sangat penting dan berguna, buku ini tidak akan membahasnya lebih jauh.

\index{Bayes' Theorem|)}
\index{tree diagram|)}
\index{conditional probability|)}
\index{probability|)}


{\input{ch_probability/TeX/conditional_probability.tex}}

\section{Pengambilan sampel dari populasi kecil}
\label{smallPop}

\noindent%
Ketika kita mengambil pengamatan dari suatu populasi sebagai sampel,
biasanya kita hanya mengambil sebagian kecil dari
individu atau kasus yang mungkin.
Namun, terkadang ukuran sampel kita cukup besar
atau populasinya cukup kecil sehingga kita
mengambil lebih dari 10\% populasi\footnote{Pedoman 10\%
  merupakan batas berdasarkan kaidah praktis untuk menentukan kapan
  pertimbangan ini menjadi lebih penting.}
\emph{tanpa
pengembalian} (artinya kita tidak mungkin
mengambil kasus yang sama dua kali).
Pengambilan bagian populasi yang cukup besar seperti ini
dapat memengaruhi cara kita menganalisis
sampel tersebut.

\begin{examplewrap}
\begin{nexample}{Dosen terkadang memilih seorang mahasiswa secara acak untuk menjawab pertanyaan. Jika setiap mahasiswa memiliki peluang yang sama untuk dipilih dan ada 15 orang di kelas Anda, berapakah peluang dosen memilih Anda untuk pertanyaan berikutnya?}
Jika ada 15 orang yang dapat ditanya dan semuanya hadir di kelas, peluangnya adalah $1/15$, atau sekitar $0.067$.
\end{nexample}
\end{examplewrap}

\begin{examplewrap}
\begin{nexample}{Jika dosen mengajukan 3 pertanyaan, berapakah peluang Anda tidak terpilih? Anggap dosen tidak memilih orang yang sama dua kali dalam satu perkuliahan.}\label{3woRep}
Untuk pertanyaan pertama, dosen memilih orang lain dengan peluang $14/15$. Ketika mengajukan pertanyaan kedua, hanya ada 14 orang yang belum ditanya. Jadi, jika Anda tidak terpilih untuk pertanyaan pertama, peluang Anda kembali tidak terpilih adalah $13/14$. Demikian pula, peluang Anda kembali tidak terpilih untuk pertanyaan ketiga adalah $12/13$, dan peluang tidak terpilih untuk ketiga pertanyaan tersebut adalah
\begin{align*}
&P(\text{tidak terpilih dalam 3 pertanyaan}) \\
&\quad = P(\text{\var{Q1}} = \text{\resp{tidak\us{}terpilih}, }\text{\var{Q2}} = \text{\resp{tidak\us{}terpilih}, }\text{\var{Q3}} = \text{\resp{tidak\us{}terpilih}.}) \\
&\quad = \frac{14}{15}\times\frac{13}{14}\times\frac{12}{13} = \frac{12}{15} = 0.80
\end{align*}
\end{nexample}
\end{examplewrap}

\begin{exercisewrap}
\begin{nexercise}
Aturan apa yang memungkinkan kita mengalikan peluang-peluang dalam Contoh~\ref{3woRep}?\footnotemark
\end{nexercise}
\end{exercisewrap}
\footnotetext{Ketiga peluang yang kita hitung sebenarnya terdiri atas satu peluang marginal, $P($\var{Q1}$ = $\resp{tidak\us{}terpilih}$)$, dan dua peluang bersyarat:
\begin{align*}
&P(\text{\var{Q2}} =  \text{\resp{tidak\us{}terpilih} }|
    \text{ \var{Q1}} = \text{\resp{tidak\us{}terpilih}}) \\
&P(\text{\var{Q3}} =  \text{\resp{tidak\us{}terpilih} }|
    \text{ \var{Q1}} = \text{\resp{tidak\us{}terpilih}, }
    \text{\var{Q2}} = \text{\resp{tidak\us{}terpilih}})
\end{align*}
Dengan Aturan Perkalian Umum, hasil kali ketiga peluang ini adalah peluang tidak terpilih dalam 3 pertanyaan.}

\D{\newpage}

\begin{examplewrap}
\begin{nexample}{Misalkan dosen memilih secara acak tanpa memperhatikan siapa yang sudah dipilih, yaitu mahasiswa dapat dipilih lebih dari sekali. Berapakah peluang Anda tidak terpilih untuk ketiga pertanyaan tersebut?}\label{3wRep}
Pemilihan-pemilihan tersebut saling independen, dan peluang tidak terpilih untuk satu pertanyaan adalah $14/15$. Jadi, kita dapat menggunakan Aturan Perkalian untuk proses yang saling independen.
\begin{align*}
&P(\text{tidak terpilih dalam 3 pertanyaan}) \\
&\quad = P(\text{\var{Q1}} = \text{\resp{tidak\us{}terpilih}, }\text{\var{Q2}} = \text{\resp{tidak\us{}terpilih}, }\text{\var{Q3}} = \text{\resp{tidak\us{}terpilih}.}) \\
&\quad = \frac{14}{15}\times\frac{14}{15}\times\frac{14}{15} = 0.813
\end{align*}
Peluang Anda untuk tidak terpilih sedikit lebih besar dibandingkan ketika dosen memilih orang yang berbeda untuk setiap pertanyaan. Namun, sekarang Anda mungkin dipilih lebih dari sekali.
\end{nexample}
\end{examplewrap}

\begin{exercisewrap}
\begin{nexercise}
Berdasarkan skenario pada Contoh~\ref{3wRep}, berapakah peluang Anda terpilih untuk menjawab ketiga pertanyaan?\footnotemark
\end{nexercise}
\end{exercisewrap}
\footnotetext{$P($terpilih untuk menjawab ketiga pertanyaan$) = \left(\frac{1}{15}\right)^3 = 0.00030$.}

Jika kita mengambil sampel dari populasi kecil \term{tanpa pengembalian}, pengamatan-pengamatan kita tidak lagi saling independen. Dalam Contoh~\ref{3woRep}, peluang tidak terpilih untuk pertanyaan kedua dikondisikan pada peristiwa bahwa Anda tidak terpilih untuk pertanyaan pertama. Dalam Contoh~\ref{3wRep}, dosen mengambil sampel mahasiswa \term{dengan pengembalian}: dosen tersebut berulang kali mengambil sampel dari seluruh kelas tanpa memedulikan siapa yang sudah dipilih. 

\begin{exercisewrap}
\begin{nexercise} \label{raffleOf30TicketsWWOReplacement}
Departemen Anda mengadakan undian berhadiah. Mereka menjual 30 tiket dan menyediakan tujuh hadiah. (a) Tiket-tiket tersebut dimasukkan ke dalam topi dan satu tiket diundi untuk setiap hadiah. Tiket diambil tanpa pengembalian, yaitu tiket yang terpilih tidak dimasukkan kembali ke dalam topi. Berapakah peluang memenangkan hadiah jika Anda membeli satu tiket? (b)~Bagaimana jika tiket diambil dengan pengembalian?\footnotemark
\end{nexercise}
\end{exercisewrap}
\footnotetext{(a) Pertama, tentukan peluang tidak menang. Tiket diambil tanpa pengembalian, yang berarti peluang Anda tidak menang pada pengundian pertama adalah $29/30$, $28/29$ pada pengundian kedua, ..., dan $23/24$ pada pengundian ketujuh. Peluang Anda tidak memenangkan hadiah adalah hasil kali peluang-peluang tersebut: $23/30$. Artinya, peluang memenangkan hadiah adalah $1 - 23/30 = 7/30 = 0.233$. (b)~Ketika tiket diambil dengan pengembalian, terdapat tujuh pengundian yang saling independen. Sekali lagi, pertama-tama kita mencari peluang tidak memenangkan hadiah: $(29/30)^7 = 0.789$. Jadi, peluang memenangkan (setidaknya) satu hadiah ketika pengambilan dilakukan dengan pengembalian adalah 0.211.}

\begin{exercisewrap}
\begin{nexercise} \label{followUpToRaffleOf30TicketsWWOReplacement}
Bandingkan jawaban Anda dalam Latihan Terbimbing~\ref{raffleOf30TicketsWWOReplacement}. Seberapa besar pengaruh metode pengambilan sampel terhadap peluang Anda memenangkan hadiah?\footnotemark
\end{nexercise}
\end{exercisewrap}
\footnotetext{Peluang memenangkan hadiah sekitar 10\% lebih besar jika pengambilan sampel dilakukan tanpa pengembalian. Namun, dengan prosedur pengambilan sampel ini, paling banyak hanya satu hadiah yang dapat dimenangkan.}

Seandainya kita mengulangi Latihan Terbimbing~\ref{raffleOf30TicketsWWOReplacement} dengan 300 tiket, bukan 30 tiket, kita akan menemukan hal yang menarik: hasilnya akan hampir identik. Peluangnya adalah 0.0233 tanpa pengembalian dan 0.0231 dengan pengembalian. Ketika ukuran sampel hanya merupakan sebagian kecil dari populasi (di bawah 10\%), pengamatan hampir saling independen meskipun pengambilan sampel dilakukan tanpa pengembalian.


{\setlength{\eoceAfterSpace}{2mm}\input{ch_probability/TeX/sampling_from_a_small_population.tex}}





%_________________
\section{Variabel acak}
\label{randomVariablesSection}

\index{random variable|(}

\noindent%
Sering kali, suatu proses berguna untuk dimodelkan dengan menggunakan apa yang
disebut \term{variabel acak}.
Model semacam ini memungkinkan kita menerapkan kerangka
matematis dan prinsip-prinsip statistika untuk
memahami dan memprediksi hasil dengan lebih baik
di dunia nyata.

\begin{examplewrap}
\begin{nexample}{Dua buku diwajibkan untuk sebuah kelas statistika: buku teks dan panduan belajar yang menyertainya. Toko buku universitas menemukan bahwa 20\% mahasiswa yang terdaftar tidak membeli satu pun dari kedua buku tersebut, 55\% hanya membeli buku teks, dan 25\% membeli keduanya; persentase ini relatif konstan dari satu semester ke semester berikutnya. Jika~ada 100 mahasiswa yang terdaftar, berapa banyak buku yang diperkirakan akan dijual toko buku tersebut kepada kelas ini?}\label{bookStoreSales}
Sekitar 20 mahasiswa tidak akan membeli satu pun buku (total 0 buku), sekitar 55 mahasiswa akan membeli satu buku (total 55 buku), dan sekitar 25 mahasiswa akan membeli dua buku (total 50 buku untuk 25 mahasiswa ini). Toko buku tersebut diperkirakan akan menjual sekitar 105 buku kepada kelas ini.
\end{nexample}
\end{examplewrap}

\begin{exercisewrap}
\begin{nexercise}
Apakah Anda akan terkejut jika toko buku tersebut menjual sedikit lebih banyak atau lebih sedikit daripada 105 buku?\footnotemark
\end{nexercise}
\end{exercisewrap}
\footnotetext{Jika toko buku menjual sedikit lebih banyak atau sedikit lebih sedikit, hal ini seharusnya tidak mengejutkan. Semoga Bab~\ref{introductionToData} telah memperjelas bahwa data yang diamati memiliki variabilitas alami. Sebagai contoh, jika kita melempar koin 100 kali, sisi kepala biasanya tidak akan muncul tepat separuh dari seluruh lemparan, tetapi jumlahnya mungkin akan mendekati separuh.}

\begin{examplewrap}
\begin{nexample}{Harga buku teks adalah \$137 dan harga panduan belajar adalah \$33. Berapa besar pendapatan yang diperkirakan akan diperoleh toko buku dari kelas yang terdiri atas 100 mahasiswa ini?}\label{bookStoreRev}
Sekitar 55 mahasiswa hanya akan membeli buku teks, sehingga menghasilkan pendapatan sebesar
\begin{align*}
\$137 \times  55 = \$7,535
\end{align*}
Sekitar 25 mahasiswa yang membeli buku teks dan panduan belajar akan membayar total sebesar
\begin{align*}
(\$137 + \$33) \times  25 = \$170 \times  25 = \$4,250
\end{align*}
Jadi, toko buku tersebut diperkirakan memperoleh sekitar $\$7,535 + \$4,250 = \$11,785$ dari 100 mahasiswa dalam satu kelas ini. Namun, mungkin terdapat \emph{variabilitas pengambilan sampel}, sehingga jumlah sebenarnya dapat sedikit berbeda.
\end{nexample}
\end{examplewrap}

\begin{figure}[h]
  \centering
  \Figure[Sebuah distribusi peluang yang tampak menyerupai histogram. Sumbu horizontal berlabel “Biaya” dan membentang dari \$0 hingga \$170. Sumbu vertikal berlabel “Peluang”. Terdapat tiga batang: batang setinggi 0.2 pada \$0, batang setinggi 0.55 pada \$137, dan batang setinggi 0.25 pada \$170. Sebuah segitiga merah ditampilkan pada rata-rata, yaitu \$117.85.]{0.6}{bookCostDist}
  \caption{Distribusi peluang pendapatan toko buku
      dari satu mahasiswa.
      Segitiga menunjukkan
      pendapatan rata-rata per mahasiswa.}
  \label{bookCostDist}
\end{figure}

\D{\newpage}

\begin{examplewrap}
\begin{nexample}{Berapakah pendapatan rata-rata per mahasiswa untuk mata kuliah ini?}\label{revFromStudent}
Total pendapatan yang diharapkan adalah \$11,785 dan terdapat 100 mahasiswa. Oleh karena itu, pendapatan yang diharapkan per mahasiswa adalah $\$11,785/100 =  \$117.85$.
\end{nexample}
\end{examplewrap}

\subsection{Nilai harapan}

\index{expectation|(}

Variabel atau proses dengan hasil numerik disebut \term{variabel acak}, dan variabel acak ini biasanya dilambangkan dengan huruf kapital seperti $X$, $Y$, atau $Z$. Jumlah uang yang dibelanjakan seorang mahasiswa untuk buku-buku statistikanya merupakan variabel acak, dan kita melambangkannya dengan $X$.

\begin{onebox}{Variabel acak}
Proses acak atau variabel yang memiliki hasil numerik.
\end{onebox}

Hasil-hasil yang mungkin dari $X$ diberi label dengan huruf kecil $x$ dan subskrip yang bersesuaian. Sebagai contoh, kita menulis $x_1=\$0$, $x_2=\$137$, dan $x_3=\$170$, yang masing-masing terjadi dengan peluang $0.20$, $0.55$, dan $0.25$. Distribusi $X$ diringkas pada Gambar~\ref{bookCostDist} dan Gambar~\ref{statSpendDist}.

\begin{figure}[h]
\centering
\begin{tabular}{l ccc r}
\hline
$i$	  & 1 & 2 & 3  & Total\\
\hline
$x_i$ & \$0 & \$137 & \$170 & --\\
$P(X=x_i)$ & 0.20 & 0.55 & 0.25 & 1.00 \\
\hline
\end{tabular}
\caption{Distribusi peluang variabel acak $X$, yang menyatakan pendapatan toko buku dari satu mahasiswa.}
\label{statSpendDist}
\end{figure}

Kita menghitung hasil rata-rata $X$ sebesar \$117.85 dalam Contoh~\ref{revFromStudent}.
Rata-rata ini disebut \term{nilai harapan} $X$, yang dilambangkan dengan $E(X)$\index{EX@$E(X)$}.
Nilai harapan suatu variabel acak dihitung dengan menjumlahkan setiap hasil yang dibobotkan menurut peluangnya:
\begin{align*}
E(X) &= 0 \times  P(X=0) + 137 \times  P(X=137) + 170 \times  P(X=170) \\
	&= 0 \times  0.20 + 137 \times  0.55 + 170 \times  0.25 = 117.85
\end{align*}

\begin{onebox}{Nilai harapan variabel acak diskret}
Jika $X$ memiliki hasil $x_1$, ..., $x_k$ dengan peluang $P(X=x_1)$, ..., $P(X=x_k)$, nilai harapan $X$ adalah jumlah setiap hasil yang dikalikan dengan peluang yang bersesuaian:
\begin{align*}
E(X)
  &= x_1 \times P(X = x_1) + \cdots + x_k\times P(X = x_k) \\
  &= \sum_{i = 1}^{k} x_i P(X = x_i)
\end{align*}
Huruf Yunani $\mu$\index{Greek!mu@mu ($\mu$)}
dapat digunakan sebagai pengganti notasi $E(X)$.
\end{onebox}

% Tata letak: pemisah halaman paksa yang usang dihapus.

Nilai harapan suatu variabel acak menyatakan hasil rata-rata. Sebagai contoh, $E(X)=117.85$ menyatakan jumlah rata-rata yang diperkirakan akan diperoleh toko buku dari satu mahasiswa, yang juga dapat kita tulis sebagai $\mu=117.85$.

Nilai harapan variabel acak kontinu juga dapat dihitung (lihat Bagian~\ref{contDist}). Namun, perhitungannya memerlukan sedikit kalkulus dan kita menundanya untuk mata kuliah lanjutan.\footnote{$\mu = \int xf(x)dx$ dengan $f(x)$ menyatakan fungsi bagi kurva kepadatan.}

Dalam fisika, nilai harapan memiliki makna yang sama dengan pusat gravitasi. Distribusi dapat direpresentasikan oleh serangkaian beban pada setiap hasil, dan rata-rata menyatakan titik keseimbangan. Hal ini ditampilkan pada Gambar~\ref{bookCostDist} dan~\ref{bookWts}. Gagasan pusat gravitasi juga berlaku untuk distribusi peluang kontinu. Gambar~\ref{contBalance} menunjukkan distribusi peluang kontinu yang seimbang di atas penyangga yang ditempatkan pada rata-rata.

\begin{figure}
\centering
\Figure[Sebuah batang digantung dengan tali, dan tiga beban tergantung pada tiga posisi berbeda di sepanjang batang. Beban-beban tersebut berada pada posisi 0, 137, dan 170 dengan bobot yang sebanding dengan peluang 0.2, 0.55, dan 0.25. Beban-beban itu seimbang karena tali yang menopang batang terletak pada rata-rata distribusi, yaitu 117.85.]{0.72}{bookWts}
\caption{Sistem beban yang merepresentasikan distribusi peluang $X$. Tali menahan distribusi pada rata-ratanya agar sistem tetap seimbang.}
\label{bookWts}
\end{figure}

\begin{figure}
\centering
\Figure[Ditampilkan sebuah distribusi yang miring ke kanan, serupa dengan tampilan histogram jika bin-binnya sangat kecil sehingga menyatu dan tampak kontinu. Distribusi ini seimbang di atas sebuah segitiga yang terletak pada rata-rata distribusi.]{0.68}{contBalance}
\caption{Distribusi kontinu juga dapat diseimbangkan pada rata-ratanya.}
\label{contBalance}
\end{figure}

\index{expectation|)}


\D{\newpage}

\subsection{Variabilitas dalam variabel acak}

Misalkan Anda mengelola toko buku universitas. Selain besarnya pendapatan yang diperkirakan akan diperoleh, Anda mungkin juga ingin mengetahui volatilitas (variabilitas) pendapatan tersebut. 

\indexthis{Varians}{variance} dan \indexthis{simpangan baku}{standard deviation} dapat digunakan untuk mendeskripsikan variabilitas suatu variabel acak. Bagian~\ref{variability}
memperkenalkan metode untuk menentukan varians dan simpangan baku suatu kumpulan data. Pertama, kita menghitung simpangan dari rata-rata ($x_i - \mu$), menguadratkan simpangan tersebut, lalu mengambil rata-ratanya untuk memperoleh varians. Untuk variabel acak, kita kembali menghitung kuadrat simpangan. Namun, kita menjumlahkannya dengan pembobotan berdasarkan peluang yang bersesuaian, seperti yang kita lakukan untuk nilai harapan. Jumlah kuadrat simpangan berbobot ini sama dengan varians, dan kita menghitung simpangan baku dengan mengambil akar kuadrat varians, seperti yang dilakukan pada Bagian~\ref{variability}.

\begin{onebox}{Rumus umum varians}
Jika $X$ memiliki hasil $x_1$, ..., $x_k$ dengan peluang $P(X=x_1)$, ..., $P(X=x_k)$ dan nilai harapan $\mu=E(X)$, varians $X$, yang dilambangkan dengan $Var(X)$ atau simbol $\sigma^2$, adalah
\begin{align*}
\sigma^2 &= (x_1-\mu)^2\times P(X=x_1) + \cdots \\
	& \qquad\quad\cdots+ (x_k-\mu)^2\times P(X=x_k) \\
	&= \sum_{j=1}^{k} (x_j - \mu)^2 P(X=x_j)
\end{align*}
Simpangan baku $X$, yang dilambangkan dengan
$\sigma$\index{Greek!sigma@sigma ($\sigma$)},
adalah akar kuadrat dari varians.
\end{onebox}

\begin{examplewrap}
\begin{nexample}{Hitung nilai harapan, varians, dan simpangan baku $X$, yaitu pendapatan toko buku dari satu mahasiswa statistika.}
Kita dapat menyusun tabel yang memuat perhitungan setiap hasil secara terpisah, lalu menjumlahkan hasil-hasilnya.
\begin{center}
\begin{tabular}{l rrr r}
\hline
$i$ & 1 & 2 & 3 & Total \\
\hline
$x_i$ & \$0 & \$137 & \$170 &  \\
$P(X=x_i)$ & 0.20 & 0.55 & 0.25 &  \\
$x_i \times  P(X=x_i)$ & 0 & 75.35 & 42.50 & 117.85 \\
\hline
\end{tabular}
\end{center}
Jadi, nilai harapannya adalah $\mu=117.85$, seperti yang telah kita hitung. Varians dapat disusun dengan memperluas tabel ini:
\begin{center}
\begin{tabular}{l rrr r}
\hline
$i$ & 1 & 2 & 3 & Total \\
\hline
$x_i$ & \$0 & \$137 & \$170 &  \\
$P(X=x_i)$ & 0.20 & 0.55 & 0.25 &  \\
$x_i \times  P(X=x_i)$ & 0 & 75.35 & 42.50 & 117.85 \\
$x_i - \mu$ & -117.85 & 19.15 & 52.15 &  \\
$(x_i-\mu)^2$ & 13888.62 &  366.72 & 2719.62 &  \\
$(x_i-\mu)^2\times P(X=x_i)$ & 2777.7 & 201.7 & 679.9 & 3659.3 \\
\hline
\end{tabular}
\end{center}
Varians $X$ adalah $\sigma^2 = 3659.3$, yang berarti simpangan bakunya adalah $\sigma = \sqrt{3659.3} = \$60.49$.
\end{nexample}
\end{examplewrap}

\begin{exercisewrap}
\begin{nexercise}
Toko buku tersebut juga menawarkan buku teks kimia seharga \$159 dan buku pendamping seharga \$41. Berdasarkan pengalaman sebelumnya, mereka mengetahui bahwa sekitar 25\% mahasiswa kimia hanya membeli buku teks, sedangkan 60\% membeli buku teks beserta buku pendamping.\footnotemark
\begin{enumerate}
\item[(a)] Berapa proporsi mahasiswa yang tidak membeli satu pun dari kedua buku tersebut? Anggap tidak ada mahasiswa yang membeli buku pendamping tanpa buku teks.
\item[(b)] Misalkan $Y$ menyatakan pendapatan dari satu mahasiswa. Tuliskan distribusi peluang $Y$, yaitu tabel yang memuat setiap hasil beserta peluangnya.
\item[(c)] Hitung pendapatan yang diharapkan dari satu mahasiswa kimia. 
\item[(d)] Tentukan simpangan baku untuk mendeskripsikan variabilitas pendapatan dari satu mahasiswa.
\end{enumerate}
\end{nexercise}
\end{exercisewrap}
\footnotetext{(a) 100\% - 25\% - 60\% = 15\% mahasiswa tidak membeli buku apa pun untuk kelas tersebut. Bagian~(b) direpresentasikan oleh dua baris pertama pada tabel di bawah ini. Nilai harapan untuk bagian~(c) diberikan sebagai total pada baris $y_i\times P(Y=y_i)$. Hasil bagian~(d) adalah akar kuadrat varians yang tercantum sebagai total pada baris terakhir: $\sigma = \sqrt{Var(Y)} = \$69.28$.
\begin{center}
\begin{tabular}{rrrrr}
  \hline
$i$ (skenario) & 1 (\resp{tanpa\us{}buku}) & 2 (\resp{buku\us{}teks}) & 3 (\resp{keduanya}) & Total \\
  \hline
$y_i$ & 0.00 & 159.00 & 200.00 &  \\
$P(Y=y_i)$ & 0.15 & 0.25 & 0.60 & \\
$y_i\times P(Y=y_i)$ & 0.00 & 39.75 & 120.00 & $E(Y) = 159.75$\\
$y_i-E(Y)$ & -159.75 & -0.75 & 40.25 & \\
$(y_i-E(Y))^2$ & 25520.06 & 0.56 & 1620.06 & \\
$(y_i-E(Y))^2\times P(Y)$ & 3828.0 & 0.1 & 972.0 & $Var(Y) \approx 4800$ \\
   \hline
\end{tabular}
\end{center}}

\subsection{Kombinasi linear variabel acak}

Sejauh ini, kita menganggap setiap variabel sebagai gambaran yang lengkap dengan sendirinya. Terkadang, kombinasi beberapa variabel lebih tepat digunakan. Sebagai contoh, waktu yang dihabiskan seseorang untuk bepergian ke tempat kerja setiap minggu dapat diuraikan menjadi beberapa perjalanan harian. Demikian pula, total keuntungan atau kerugian dalam suatu portofolio saham merupakan jumlah keuntungan dan kerugian dari komponen-komponennya.

\begin{examplewrap}
\begin{nexample}{John pergi ke tempat kerja lima hari dalam seminggu. Kita akan menggunakan $X_1$ untuk menyatakan waktu perjalanannya pada hari Senin, $X_2$ untuk menyatakan waktu perjalanannya pada hari Selasa, dan seterusnya. Tuliskan persamaan menggunakan $X_1$, ..., $X_5$ yang menyatakan waktu perjalanannya selama seminggu, yang dilambangkan dengan $W$.}
Total waktu perjalanan mingguannya adalah jumlah kelima nilai harian:
\begin{align*}
W = X_1 + X_2 + X_3 + X_4 + X_5
\end{align*}
Menguraikan waktu perjalanan mingguan $W$ menjadi beberapa bagian menyediakan kerangka untuk memahami setiap sumber keacakan dan berguna dalam memodelkan $W$.
\end{nexample}
\end{examplewrap}

\D{\newpage}

\begin{examplewrap}
\begin{nexample}{Waktu rata-rata yang dibutuhkan John untuk pergi ke tempat kerja setiap hari adalah 18 menit. Berapa waktu perjalanan rata-rata yang Anda harapkan selama seminggu?}
Kita diberi tahu bahwa rata-rata (yaitu nilai harapan) waktu perjalanan adalah 18 menit per hari: $E(X_i) = 18$. Untuk memperoleh waktu yang diharapkan bagi jumlah lima hari tersebut, kita dapat menjumlahkan waktu yang diharapkan untuk setiap hari:
\begin{align*}
E(W) &= E(X_1 + X_2 + X_3 + X_4 + X_5) \\
	&= E(X_1) + E(X_2) + E(X_3) + E(X_4) + E(X_5) \\
	&= 18 + 18 + 18 + 18 + 18 = 90\text{ menit}
\end{align*}
Nilai harapan waktu total sama dengan jumlah nilai harapan waktu masing-masing. Secara lebih umum, nilai harapan jumlah beberapa variabel acak selalu sama dengan jumlah nilai harapan setiap variabel acak.
\end{nexample}
\end{examplewrap}

\begin{exercisewrap}
\begin{nexercise}
\label{elenaIsSellingATVAndBuyingAToasterOvenAtAnAuction}%
Elena menjual sebuah TV dalam lelang tunai dan juga bermaksud membeli oven pemanggang dalam lelang tersebut. Jika $X$ menyatakan keuntungan dari penjualan TV dan $Y$ menyatakan biaya oven pemanggang, tuliskan persamaan yang menyatakan perubahan bersih uang tunai Elena.\footnotemark
\end{nexercise}
\end{exercisewrap}
\footnotetext{Elena akan memperoleh $X$ dolar dari TV, tetapi membelanjakan $Y$ dolar untuk oven pemanggang: $X-Y$.}

\begin{exercisewrap}
\begin{nexercise}
Berdasarkan lelang sebelumnya, Elena memperkirakan akan memperoleh sekitar \$175 dari TV dan membayar sekitar \$23 untuk oven pemanggang. Secara keseluruhan, berapa banyak uang yang diperkirakan akan ia peroleh atau belanjakan?\footnotemark
\end{nexercise}
\end{exercisewrap}
\footnotetext{$E(X-Y) = E(X) - E(Y) = 175 - 23 = \$152$. Elena diperkirakan akan memperoleh sekitar \$152.}

\begin{exercisewrap}
\begin{nexercise} \label{explainWhyThereIsUncertaintyInTheSum}
Apakah Anda akan terkejut jika waktu perjalanan mingguan John tidak tepat 90 menit atau jika Elena tidak memperoleh tepat \$152? Jelaskan.\footnotemark
\end{nexercise}
\end{exercisewrap}
\footnotetext{Tidak, karena kemungkinan terdapat variabilitas. Sebagai contoh, kondisi lalu lintas berubah dari satu hari ke hari berikutnya, dan harga lelang akan bervariasi bergantung pada kualitas barang serta minat para peserta.}

Dua konsep penting mengenai kombinasi variabel acak telah diperkenalkan. Pertama, suatu nilai akhir terkadang dapat dinyatakan dalam persamaan sebagai jumlah bagian-bagiannya. Kedua, secara intuitif, memasukkan nilai rata-rata setiap bagian ke dalam persamaan ini akan memberikan nilai total rata-rata yang kita harapkan. Poin kedua ini perlu diperjelas -- hal tersebut dijamin benar untuk apa yang disebut \emph{kombinasi linear variabel acak}.

\term{Kombinasi linear} dua variabel acak $X$ dan $Y$ adalah istilah untuk menyatakan kombinasi
\begin{align*}
aX + bY
\end{align*}
dengan $a$ dan $b$ sebagai bilangan tetap yang diketahui. Untuk waktu perjalanan John, terdapat lima variabel acak -- satu untuk setiap hari kerja -- dan setiap variabel acak dapat ditulis dengan koefisien tetap sebesar 1:
\begin{align*}
1X_1 + 1 X_2 + 1 X_3 + 1 X_4 + 1 X_5
\end{align*}
Untuk keuntungan atau kerugian bersih Elena, variabel acak $X$ memiliki
koefisien +1 dan variabel acak $Y$ memiliki
koefisien~-1.

\D{\newpage}

Ketika mempertimbangkan rata-rata kombinasi linear variabel acak, kita dapat memasukkan rata-rata setiap variabel acak lalu menghitung hasil akhirnya. Beberapa contoh kombinasi nonlinear variabel acak -- yaitu kasus ketika kita tidak dapat sekadar memasukkan rata-rata -- diberikan dalam catatan kaki.\footnote{Jika $X$ dan $Y$ merupakan variabel acak, pertimbangkan kombinasi berikut: $X^{1+Y}$, $X\times Y$, $X/Y$. Dalam kasus seperti ini, memasukkan nilai rata-rata setiap variabel acak dan menghitung hasilnya umumnya tidak akan menghasilkan nilai rata-rata yang akurat bagi hasil akhir.}

\begin{onebox}{Kombinasi linear variabel acak dan hasil rata-rata}
Jika $X$ dan $Y$ merupakan variabel acak, kombinasi linear kedua variabel acak tersebut diberikan oleh
\begin{align*}
aX + bY
\end{align*}
dengan $a$ dan $b$ sebagai bilangan tetap. Untuk menghitung nilai rata-rata kombinasi linear variabel acak, masukkan rata-rata setiap variabel acak lalu hitung hasilnya:
\begin{align*}
a\times E(X) + b\times E(Y)
\end{align*}
Ingat bahwa nilai harapan sama dengan rata-rata, misalnya $E(X) = \mu_X$.
\end{onebox}

\begin{examplewrap}
\begin{nexample}{Leonard telah menginvestasikan \$6000 di Caterpillar Inc
    (kode saham: CAT) dan \$2000 di Exxon Mobil Corp (XOM).
    Jika $X$ menyatakan perubahan saham Caterpillar pada bulan depan
    dan $Y$ menyatakan perubahan saham Exxon Mobil
    pada bulan depan, tuliskan persamaan yang mendeskripsikan berapa banyak
    uang yang akan diperoleh atau hilang dari saham Leonard selama
    bulan tersebut.}
  Untuk menyederhanakan, kita menganggap $X$ dan $Y$ tidak dinyatakan
  dalam persen, melainkan dalam bentuk desimal
  (misalnya, jika saham Caterpillar naik 1\%, maka $X=0.01$;
  atau jika turun 1\%, maka $X=-0.01$).
  Dengan demikian, kita dapat menuliskan persamaan keuntungan Leonard sebagai
  \begin{align*}
  \$6000\times X + \$2000\times Y
  \end{align*}
  Jika kita memasukkan perubahan nilai saham untuk $X$ dan $Y$,
  persamaan ini memberikan perubahan nilai portofolio saham Leonard
  selama bulan tersebut. Nilai positif menyatakan keuntungan,
  sedangkan nilai negatif menyatakan kerugian.
\end{nexample}
\end{examplewrap}

\begin{exercisewrap}
\begin{nexercise}\label{expectedChangeInLeonardsStockPortfolio}
Saham Caterpillar akhir-akhir ini naik
sebesar 2.0\% dan saham Exxon Mobil sebesar 0.2\% per bulan.
Hitung perubahan yang diharapkan pada portofolio saham Leonard
untuk bulan depan.\footnotemark
\end{nexercise}
\end{exercisewrap}
% library(openintro); d <- stocks_18; cols <- 2:4; apply(d[, cols], 2, mean); apply(d[, cols], 2, sd)
\footnotetext{%
  $E(\$6000\times X + \$2000\times Y) =
    \$6000\times 0.020 + \$2000\times 0.002 = \$124$.}

\begin{exercisewrap}
\begin{nexercise}
Anda seharusnya menemukan bahwa Leonard mengharapkan keuntungan positif
dalam Latihan Terbimbing~\ref{expectedChangeInLeonardsStockPortfolio}.
Namun, apakah Anda akan terkejut jika ia sebenarnya mengalami
kerugian pada bulan ini?\footnotemark
\end{nexercise}
\end{exercisewrap}
\footnotetext{Tidak.
  Walaupun saham cenderung naik dari waktu ke waktu, nilainya sering
  bergejolak dalam jangka pendek.}


\D{\newpage}

\subsection{Variabilitas dalam kombinasi linear variabel acak}
\label{var_lin_combo_of_RVs}

Mengukur hasil rata-rata dari kombinasi linear
variabel acak memang berguna, tetapi kita juga perlu
memahami ketidakpastian yang berkaitan dengan
hasil total kombinasi variabel acak tersebut.
Keuntungan atau kerugian bersih yang diharapkan dari portofolio saham Leonard
telah dipertimbangkan dalam Latihan Terbimbing~\ref{expectedChangeInLeonardsStockPortfolio}.
Namun, volatilitas portofolio ini belum
dibahas secara kuantitatif.
Sebagai contoh, meskipun keuntungan bulanan rata-rata menurut data mungkin
sekitar \$124, keuntungan tersebut tidak terjamin.
Gambar~\ref{changeInLeonardsStockPortfolioFor36Months}
menunjukkan perubahan bulanan pada portofolio seperti milik Leonard selama
periode tiga tahun.
Keuntungan dan kerugiannya sangat bervariasi, dan mengukur
fluktuasi ini penting ketika berinvestasi dalam saham.

\begin{figure}[ht]
\centering
\Figure[Diagram titik ditumpangkan pada diagram kotak untuk variabel “Imbal Hasil Bulanan Selama 3 Tahun”. Bagian kotak membentang kira-kira dari -200 hingga 450, dengan garis median sekitar 200. Garis kumis membentang ke ujung bawah sekitar -600 dan ujung atas sekitar 1050. Terdapat satu titik di luar garis kumis bawah pada sekitar -1400. Pada diagram titik, titik-titik tersebar cukup merata di sepanjang bagian kotak dan garis kumis, dengan satu pengecualian berupa titik pada -1400.]{0.6}{changeInLeonardsStockPortfolioFor36Months}
\caption{Perubahan portofolio seperti milik Leonard selama 36 bulan,
    dengan \$6000 diinvestasikan dalam saham Caterpillar dan \$2000 dalam
    saham Exxon Mobil.}
\label{changeInLeonardsStockPortfolioFor36Months}
\end{figure}

Seperti dalam banyak kasus sebelumnya,
kita menggunakan varians dan simpangan baku untuk mendeskripsikan
ketidakpastian yang berkaitan dengan imbal hasil bulanan Leonard.
Untuk itu, varians imbal hasil bulanan setiap saham
akan berguna, dan nilainya ditampilkan pada
Gambar~\ref{sumStatOfCATXOM}.
Imbal hasil kedua saham tersebut hampir saling independen.

\begin{figure}
\centering
\begin{tabular}{lrrr}
\hline
    & Rata-rata ($\bar{x}$) & Simpangan baku ($s$) &
	    Varians ($s^2$) \\
\hline
CAT & 0.0204  & 0.0757  & 0.0057 \\
XOM & 0.0025  & 0.0455  & 0.0021 \\
\hline
\end{tabular}
\caption{Rata-rata, simpangan baku, dan varians saham
    CAT dan XOM.
    Statistik ini diestimasi dari data saham
    historis, sehingga digunakan notasi untuk
    statistik sampel.}
\label{sumStatOfCATXOM}
\end{figure}

Di sini kita menggunakan persamaan dari teori probabilitas untuk
mendeskripsikan ketidakpastian imbal hasil bulanan Leonard;
pembuktian metode ini diserahkan kepada mata kuliah
khusus probabilitas.
Varians kombinasi linear variabel acak
dapat dihitung dengan memasukkan varians setiap
variabel acak dan menguadratkan koefisien
variabel-variabel acak tersebut:
\begin{align*}
Var(aX + bY) = a^2\times Var(X) + b^2\times Var(Y)
\end{align*}
Perlu diperhatikan bahwa kesamaan ini mengasumsikan
variabel-variabel acak tersebut saling independen;
%\Comment{new description here about if independence is broken}
jika independensi tidak terpenuhi, persamaan ini
perlu dimodifikasi. Topik tersebut akan
dibahas dalam mata kuliah selanjutnya.
Persamaan ini dapat digunakan untuk menghitung varians
imbal hasil bulanan Leonard:
\begin{align*}
Var(6000\times X + 2000\times Y)
	&= 6000^2\times Var(X) + 2000^2\times Var(Y) \\
	&= 36,000,000\times 0.0057 + 4,000,000\times 0.0021 \\
	&\approx 213,600
% sum(c(36e6, 4e6) * c(0.0057, 0.0021))
\end{align*}
Simpangan baku dihitung sebagai akar kuadrat
varians: $\sqrt{213,600} = \$463$.
Meskipun imbal hasil bulanan rata-rata sebesar \$124 dari
investasi \$8000 bukanlah jumlah yang kecil,
imbal hasil bulanan tersebut begitu bergejolak sehingga Leonard sebaiknya
tidak mengharapkan pendapatan ini sangat stabil.

\begin{onebox}{Variabilitas kombinasi linear variabel acak}
Varians kombinasi linear variabel acak dapat dihitung dengan menguadratkan konstanta, mengganti variabel acak dengan variansnya, lalu menghitung hasilnya:
\begin{align*}
Var(aX + bY) = a^2\times Var(X) + b^2\times Var(Y)
\end{align*}
Persamaan ini berlaku selama variabel-variabel acak tersebut saling independen. Simpangan baku kombinasi linear dapat diperoleh dengan mengambil akar kuadrat varians.
\end{onebox}

\begin{examplewrap}
\begin{nexample}{Misalkan waktu perjalanan harian John memiliki simpangan baku 4 menit. Berapakah ketidakpastian total waktu perjalanannya selama seminggu?} \label{sdOfJohnsCommuteWeeklyTime}
Ekspresi untuk waktu perjalanan John adalah
\begin{align*}
X_1 + X_2 + X_3 + X_4 + X_5
\end{align*}
Setiap koefisien bernilai 1, dan varians waktu perjalanan setiap hari adalah $4^2=16$. Jadi, varians total waktu perjalanan mingguan adalah
\begin{align*}
&\text{varians }= 1^2 \times  16 + 1^2 \times  16 + 1^2 \times  16 + 1^2 \times  16 + 1^2 \times  16 = 5\times 16 = 80 \\
&\text{simpangan baku } = \sqrt{\text{varians}} = \sqrt{80} = 8.94
\end{align*}
Simpangan baku waktu perjalanan mingguan John ke tempat kerja adalah sekitar 9 menit.
\end{nexample}
\end{examplewrap}

\begin{exercisewrap}
\begin{nexercise}
Perhitungan dalam Contoh~\ref{sdOfJohnsCommuteWeeklyTime} bergantung pada asumsi penting: waktu perjalanan setiap hari saling independen dari waktu perjalanan pada hari-hari lain dalam minggu tersebut. Menurut Anda, apakah asumsi ini valid? Jelaskan.\footnotemark
\end{nexercise}
\end{exercisewrap}
\footnotetext{Salah satu hal yang perlu dipertimbangkan adalah apakah pola lalu lintas cenderung memiliki siklus mingguan (misalnya, hari Jumat mungkin lebih buruk daripada hari lainnya). Jika demikian dan John berkendara, asumsi tersebut mungkin tidak masuk akal. Namun, jika John berjalan kaki ke tempat kerja, waktu perjalanannya mungkin tidak dipengaruhi oleh siklus lalu lintas mingguan.}

\begin{exercisewrap}
\begin{nexercise}\label{elenaIsSellingATVAndBuyingAToasterOvenAtAnAuctionVariability}
Pertimbangkan kembali dua lelang Elena dari Latihan Terbimbing~\ref{elenaIsSellingATVAndBuyingAToasterOvenAtAnAuction} pada halaman~\pageref{elenaIsSellingATVAndBuyingAToasterOvenAtAnAuction}. Misalkan kedua lelang ini kira-kira saling independen dan variabilitas harga lelang TV serta oven pemanggang dapat dideskripsikan menggunakan simpangan baku \$25 dan \$8. Hitung simpangan baku keuntungan bersih Elena.\footnotemark
\end{nexercise}
\end{exercisewrap}
\footnotetext{Persamaan untuk Elena dapat ditulis sebagai
\begin{align*}
(1)\times X + (-1)\times Y
\end{align*}
Varians $X$ dan $Y$ adalah 625 dan 64. Kita menguadratkan koefisien dan memasukkan nilai varians:
\begin{align*}
(1)^2\times Var(X) + (-1)^2\times Var(Y) = 1\times 625 + 1\times 64 = 689
\end{align*}
Varians kombinasi linear tersebut adalah 689, dan simpangan bakunya adalah akar kuadrat 689: sekitar \$26.25.}

Pertimbangkan kembali Latihan Terbimbing~\ref{elenaIsSellingATVAndBuyingAToasterOvenAtAnAuctionVariability}. Koefisien negatif untuk $Y$ dalam kombinasi linear hilang ketika kita menguadratkan koefisien. Hal ini berlaku secara umum: tanda negatif dalam kombinasi linear tidak memengaruhi variabilitas yang dihitung untuk kombinasi linear, tetapi memengaruhi perhitungan nilai harapan.

\index{random variable|)}


{\input{ch_probability/TeX/random_variables.tex}}





%_________________
\section{Continuous distributions}
\label{contDist}

\noindent%
So far in this chapter we've discussed cases
where the outcome of a variable is discrete.
In this section, we consider a context where
the outcome is a continuous numerical variable.

\index{data!US adult heights|(}
\index{hollow histogram|(}
\begin{examplewrap}
\begin{nexample}{Figure~\ref{fdicHistograms} shows a few
    different hollow histograms for the heights of US adults.
    How does changing the number of bins allow you to make
    different interpretations of the data?}
  \label{usHeights}%
  Adding more bins provides greater detail.
  This sample is extremely large, which is why much smaller
  bins still work well.
  Usually we do not use so many bins with smaller sample
  sizes since small counts per bin mean the bin heights
  are very volatile.
\end{nexample}
\end{examplewrap}

\begin{figure}[ht]
  \centering
  \Figure[Four hollow histograms are shown for the US adult heights in centimeters with varying bin widths. The difference in appears will first be discussed, and then the shape of the last, most detailed histogram will be given. The first histogram has about 6 bins with values that appear to be non-zero, so the outline is very boxy. The second histogram has about 12 non-zero bins, and appears a bit more continuous and less boxy than the first histogram. The third histogram has about 25 non-zero bins, and the hollow histogram outline is starting to look somewhat smoother. The last histogram has about 50 bins, and due to the large number of bins, the distribution looks quite smooth, in that no steps from one bin to the next is a substantial jump or drop in height. Next, this last histogram is described: The bin heights are about zero until 147, then they steadily climb up to about 155 before steeply climbing a little until 157 and then steadily climb to a peak at about 165. From here the histogram declines about 10\% from its peak at 170, at which point the decline is more gradual until about 183, at which point it descends rapidly until about 187 where it begins to descend more slowly as it approaches 0. At about 200, the bin heights have essentially hit zero and stay there.]{}{fdicHistograms}
  \caption{Four hollow histograms of US adults heights
      with varying bin widths.}
  \label{fdicHistograms}
\end{figure}

\begin{examplewrap}
\begin{nexample}{What proportion of the sample is between \resp{180} cm and \resp{185} cm tall (about 5'11" to 6'1")?}\label{contDistProb}
We can add up the heights of the bins in the range \resp{180} cm and \resp{185} and divide by the sample size. For instance, this can be done with the two shaded bins shown in Figure~\ref{usHeightsHist180185}. The two bins in this region have counts of 195,307 and 156,239 people, resulting in the following estimate of the probability:
\begin{align*}
\frac{195307 + 156239}{\text{3,000,000}} = 0.1172
\end{align*}
This fraction is the same as the proportion of the histogram's area that falls in the range \resp{180} to \resp{185} cm.
\end{nexample}
\end{examplewrap}

\begin{figure}[h]
  \centering
  \Figure[A histogram for heights is shown, with the two histogram bins between 180 and 185 centimeters are shaded, representing the individuals with heights between 180 and 185 centimeters.]{0.9}{usHeightsHist180185}
  \caption{A histogram with bin sizes of 2.5 cm.
      The shaded region represents individuals with
      heights between \resp{180} and \resp{185} cm.}
  \label{usHeightsHist180185}
\end{figure}


\D{\newpage}

\subsection{From histograms to continuous distributions}

Examine the transition from a boxy hollow histogram in the top-left of Figure~\ref{fdicHistograms} to the much smoother plot in the lower-right. In this last plot, the bins are so slim that the hollow histogram is starting to resemble a smooth curve. This suggests the population height as a \emph{continuous} numerical variable might best be explained by a curve that represents the outline of extremely slim bins.

This smooth curve represents a
\termsub{probability density function}
    {probability!density function}
(also called a \term{density} or \term{distribution}), and such a curve is shown in Figure~\ref{fdicHeightContDist} overlaid on a histogram of the sample. A density has a special property: the total area under the density's curve is 1. 

\begin{figure}[tbh]
\centering
\Figure[A histogram for heights of US adults is shown with an overlaid continuous line along the heights of the bins. This continuous line is smooth and would represent what a hollow histogram would look like if we had infinite data and the bin width was so small that the boxy outline of the histogram looks continuous. This is called a "continuous probability density".]{0.9}{fdicHeightContDist}
\caption{The continuous probability distribution of heights for US adults.}
\label{fdicHeightContDist}
\end{figure}

\index{hollow histogram|)}


\D{\newpage}

\subsection{Probabilities from continuous distributions}

We computed the proportion of individuals with heights \resp{180} to \resp{185} cm in Example~\ref{contDistProb} as a fraction:
\begin{align*}
\frac{\text{number of people between \resp{180} and \resp{185}}}{\text{total sample size}}
\end{align*}
We found the number of people with heights between \resp{180} and \resp{185} cm by determining the fraction of the histogram's area in this region. Similarly, we can use the area in the shaded region under the curve to find a probability (with the help of a computer):
\begin{align*}
P(\text{\var{height} between \resp{180} and \resp{185}})
	= \text{area between \resp{180} and \resp{185}}
	= 0.1157
\end{align*}
The probability that a randomly selected person is between \resp{180} and \resp{185} cm is 0.1157. This is very close to the estimate from Example~\ref{contDistProb}: 0.1172. 

\begin{figure}[h]
  \centering
  \Figure[A density curve for heights is shown, with the region between 180 and 185 centimeters being shaded.]{0.7}{fdicHeightContDistFilled}
  \caption{Density for heights in the US adult population
      with the area between 180 and 185 cm shaded.
      Compare this plot with Figure~\ref{usHeightsHist180185}.}
  \label{fdicHeightContDistFilled}
\end{figure}

\begin{exercisewrap}
\begin{nexercise}
Three US adults are randomly selected. The probability a single adult is between \resp{180} and \resp{185} cm is 0.1157.\footnotemark\vspace{-1.5mm}
\begin{enumerate}
\setlength{\itemsep}{0mm}
\item[(a)] What is the probability that all three are between \resp{180} and \resp{185} cm tall?
\item[(b)] What is the probability that none are between \resp{180} and \resp{185} cm?
\end{enumerate}
\end{nexercise}
\end{exercisewrap}
\footnotetext{Brief answers:
  (a) $0.1157 \times 0.1157 \times 0.1157 = 0.0015$.
  (b) $(1-0.1157)^3 = 0.692$}

\begin{examplewrap}
\begin{nexample}{What is the probability that a randomly selected person is \textbf{exactly} \resp{180}~cm? Assume you can measure perfectly.}
\label{probabilityOfExactly180cm}
This probability is zero. A person might be close to \resp{180} cm, but not exactly \resp{180} cm tall. This also makes sense with the definition of probability as area; there is no area captured between \resp{180}~cm and \resp{180}~cm.
\end{nexample}
\end{examplewrap}

\begin{exercisewrap}
\begin{nexercise}
Suppose a person's height is rounded to the nearest centimeter. Is there a chance that a random person's \textbf{measured} height will be \resp{180} cm?\footnotemark
\end{nexercise}
\end{exercisewrap}
\footnotetext{This has positive probability. Anyone between \resp{179.5} cm and \resp{180.5} cm will have a \emph{measured} height of \resp{180} cm. This is probably a more realistic scenario to encounter in practice versus Example~\ref{probabilityOfExactly180cm}.}

\index{data!US adult heights|)}


{\input{ch_probability/TeX/continuous_distributions.tex}}
